% v4.9.0 post-release layers
% Appended layers v4.9.0-r01 to r16, carried verbatim from the full-monolith working source
% Quantum_Sphaera_Companion_v4_8_0_source.tex (SHA256 613d9ecb107dccc79f5b686941a352635422e83f378b8f4234aaa537a85352de),
% lines 67718-70131.  Not a preserved v4.7 fragment; earlier layers are not rewritten.
% Content SHA256 of the carried block: ed05bb7ef631a4ae91b91b8789806d5c4ee43a053e5c7266af13d48e31abfb7b

% ============================================================================
\clearpage
\section*{v4.9.0-r01 Source-Only: The Photon Frame as Neutral Axis and the Conversion of Constants}
\addcontentsline{toc}{section}{v4.9.0-r01 Source-Only: The Photon Frame as Neutral Axis and the Conversion of Constants}
\label{sec:v490-r01-neutral-axis-conversion}
% ============================================================================

\paragraph{Purpose.}
Finding~(a) of \S\ref{subsec:v480-r18-findings} records that the conversion of constants between the matter and the
photon frame uses \(\gL=\cosh2\eta\), the matter--tachyon cross term, where the matter--photon pair is a neighbour pair
with factor \(\cosh\eta\).  This layer records where the factor came from, the working assumption under which it was
chosen, and what the assumption implies for a conversion law over three frames.  The computations are
\texttt{llm/k37\_gl\_umrechnungen.py}.  No value is compared with data.  Earlier layers are not rewritten.

\noindent\textbf{Source-only status.}
This v4.9.0-r01 layer is source-only.  It does not build a PDF, does not assign a new DOI,
and does not authorize a public release.  The v4.9.0 Version DOI
\href{https://doi.org/10.5281/zenodo.23534954}{10.5281/zenodo.23534954} and the Concept DOI
\href{https://doi.org/10.5281/zenodo.19210728}{10.5281/zenodo.19210728} are unchanged by it.

\subsection*{v4.9.0-r01 provenance of the factor}
\label{subsec:v490-r01-provenance}

In the archived versions v2.28.0 to v2.29.0 the splitting factor reads
``\(\gL=\gBI/\gstring\approx1.197\)'', while the same versions already carry \(1/\bsym\approx1.820\), that is
\(\bsym\approx0.5495\).  The value \(1.197\) is \((1-\bsym^2)^{-1/2}\), the Lorentz factor of the matter-frame
offset; the ratio \(\gBI/\gstring\) is \(1.864\).  In v2.29.1 the value was changed to \(1.864\) and the change recorded
as an arithmetic correction (``previous version used incorrect \(\gL\approx1.197\)''); the note carried in
Corollary~\ref{cor:gue-frame} calls the earlier value \(\gBI\), which it is not (\(\gBI\approx0.2375\)).  Since then
every conversion uses \(1.864\).

\paragraph{Author's statement (2026-10-21).}
The photon frame was assumed to be neutral, and the work was meant to proceed on that neutral axis first.  The archive
records the change of value only as an arithmetic correction and does not mention a neutral axis.

\paragraph{Theoremlet v4.9.0-r01.1.}
The two values the factor has carried are the two cross terms of Theoremlet r18.1:
\[
1.197=\cosh\eta=\sqrt{\tfrac{1+\gL}{2}}\quad\text{(matter--photon)},\qquad
1.864=\cosh2\eta=\gL\quad\text{(matter--tachyon)},
\]
with \(\eta=\operatorname{artanh}\bsym\).  Both follow from the single input \(\gBI/\gstring\).

\paragraph{Proof sketch.}
\(\cosh2\eta=2\cosh^2\eta-1\) and \(\cosh\eta=(1-\tanh^2\eta)^{-1/2}\).  \(\square\)

\subsection*{v4.9.0-r01 the neutral assumption in Reading~B}
\label{subsec:v490-r01-neutral}

The assumption has structural support in the complex-vector reading.  The photon is the sector fixed by the whole duality
circle (\S\ref{subsec:v480-r07-duality-angle}), the branch point of the duality holonomy (Theoremlet r19.2), and the
centre of the real \(s\)-axis on which matter and tachyon are mirror images (\S\ref{subsec:v480-r19-photon-centred}).
On the rapidity axis of Theorem~\ref{thm:symmetric-frame} it sits at rapidity \(0\), between matter at \(+\eta\) and
tachyon at \(-\eta\).  If the photon frame is neutral, a value on the neutral axis is reached from the matter frame by the
neighbour step, and the factor carried before v2.29.1 is the factor of that step.

\paragraph{Theoremlet v4.9.0-r01.2.}
A conversion law \(X^{(b)}=\gamma_{ab}^{-k}X^{(a)}\) with \(\gamma_{ab}\) the Lorentz factor of the pair is not
consistent over the three frames: from matter to tachyon through the photon it gives
\((\cosh\eta)^{2k}=\bigl(\tfrac{1+\gL}{2}\bigr)^k\), directly it gives \(\gL^{\,k}\) (for \(k=1\): \(1.432\) against
\(1.864\); for \(k=6\): \(8.63\) against \(42.0\)).  A law with the Doppler factor, \(X^{(b)}=e^{-k\eta_{ab}}X^{(a)}\),
composes exactly.

\paragraph{Proof sketch.}
\(\cosh^2\eta\neq\cosh2\eta\) for \(\eta\neq0\), while \(e^{\eta}e^{\eta}=e^{2\eta}\).  \(\square\)

\paragraph{Consequence.}
With the Lorentz-factor law one pair at a time is consistent.  Working on one axis first, as the author describes, is the
regime in which the law of \S\ref{sec:phase-constants} is free of contradiction; the obstruction of r01.2 appears only
once both pairs are used.  Under the neutral assumption the primary pair is matter--photon, with factor \(\cosh\eta\);
\(\gL\) is the factor across the neutral axis, from matter to tachyon.

\subsection*{v4.9.0-r01 a frame-order finding in the Hawking proposition}
\label{subsec:v490-r01-hawking}

Items~(i) and~(ii) of Proposition~\ref{prop:hawking-gBI-tachyon} give \(T_H^{(\mathrm{photon})}=\gL\,T_H^{(\mathrm{matter})}\)
and \(T_H^{(\mathrm{tachyon})}=\gL^{-1}T_H^{(\mathrm{matter})}\).  Reciprocal factors place photon and tachyon on opposite
sides of matter.  On the axis of Theorem~\ref{thm:symmetric-frame} the photon lies between matter and tachyon.  A law along
the axis moves the tachyon value beyond the photon value in the same direction; a mirror law about a neutral photon gives
the tachyon the matter value.  Neither gives reciprocal factors, so item~(ii) carries a frame order that does not match the
axis.  This is also the one statement of the main body in which the same factor serves both pairs.

\subsection*{v4.9.0-r01 ways forward, recorded and not chosen}
\label{subsec:v490-r01-ways}

\begin{enumerate}[leftmargin=2em]
\item[(A)] Keep \(\gL\) and the Lorentz-factor law, and read each conversion of \S\ref{sec:phase-constants} as one across
the neutral axis (matter--tachyon).  Statements whose endpoint is the photon frame then need the neighbour factor.
\item[(B)] Use the neighbour factor \(\cosh\eta=1.197\) for matter--photon, the value carried before v2.29.1, in line with
the neutral assumption.  This changes the photon-frame numbers of the main body (table below).
\item[(C)] Use a composable law with the Doppler factor, \(e^{\eta}=\sqrt{(1+\bsym)/(1-\bsym)}=1.854\) per neighbour step.
It is consistent over three frames.  At \(\bsym\) its neighbour factor lies \(0.55\%\) below \(\gL\), so the numbers of the
main body change little while their meaning changes.  By Theoremlet r01.3 below the closeness is structural: in the
continuum limit of Theorem~\ref{thm:tribonacci} the two factors are equal.
\end{enumerate}

\paragraph{Theoremlet v4.9.0-r01.3.}
The full-span Lorentz factor equals the neighbour-step Doppler factor, \(\cosh2\eta=e^{\eta}\), exactly when \(e^{\eta}=1\)
or \(e^{\eta}=\mathbb T\), the tribonacci constant of Theorem~\ref{thm:tribonacci}.  At that point
\(\cosh2\eta=e^{\eta}=\coth\eta=\mathbb T=1.83929\), \(\beta=0.54369\): the self-referential condition
\(\gL=1/\bsym(\gL)\) of Theorem~\ref{thm:tribonacci} is the condition that the matter--tachyon Lorentz factor and the
matter--photon Doppler factor coincide.

\paragraph{Proof sketch.}
With \(x=e^{\eta}\), \(\cosh2\eta=(x^2+x^{-2})/2\), so \(\cosh2\eta=x\) is \(x^4-2x^3+1=0\), which factors as
\((x-1)(x^3-x^2-x-1)=0\).  Further \(\coth\eta=(x^2+1)/(x^2-1)\), and \(\coth\eta=x\) is again \(x^3-x^2-x-1=0\).
Checked numerically to \(10^{-16}\).  \(\square\)

\paragraph{Consequence for the neutral axis.}
In the continuum limit the author's choice of \(\gL\) for the step from matter to the neutral photon frame is exactly the
factor a composable law assigns to that step: reading (A) and reading (C) agree there.  With the measured
\(\gL=\gBI/\gstring\), which lies \(1.37\%\) above \(\mathbb T\) (the lattice correction of
Remark~\ref{rem:op1-tribonacci}), they differ by \(0.55\%\).  The Lorentz-factor law of reading (B) does not share this
property.

\begin{center}
\small
\begin{tabular}{@{}lrrr@{}}
\toprule
quantity in the main body & \(\gL\) & neighbour, Lorentz & neighbour, Doppler\\
\midrule
\(\lP^{(\mathrm{matter})}/\lP^{(\mathrm{photon})}\) (Theorem~\ref{thm:planck-photon}), power \(3/2\) & \(2.546\) & \(1.309\) & \(2.525\)\\
vacuum-energy ratio, power \(6\) & \(42.0\) & \(2.94\) & \(40.6\)\\
\(m_\nu^{(\mathrm{matter})}/m_\nu^{(\mathrm{photon})}\) (Proposition~\ref{prop:k-neutrino}), power \(4\) & \(12.08\) & \(2.05\) & \(11.82\)\\
matter error (Corollary~\ref{cor:matter-error}), power \(5\) & \(22.53\) & \(2.46\) & \(21.92\)\\
\bottomrule
\end{tabular}
\end{center}
The table is arithmetic on the formulas of the main body; photon-frame values are not measured.  Statements that are
compared with data, among them Corollary~\ref{cor:falsifiable} and Proposition~\ref{prop:muon-g2}, are not re-evaluated
here; a re-evaluation under (A), (B) or (C) is a prediction and requires a pre-registration note.

\subsection*{v4.9.0-r01 register entries}
\label{subsec:v490-r01-register}

\begin{center}
\small
\begin{tabular}{@{}>{\raggedright\arraybackslash}p{0.34\textwidth}>{\raggedright\arraybackslash}p{0.58\textwidth}@{}}
\toprule
statement & status after v4.9.0-r01\\
\midrule
conversion law of \S\ref{sec:phase-constants} & factor \(\gL\) is the matter--tachyon term; consistent for one pair at a time (r01.2); under the neutral photon assumption the matter--photon factor is \(\cosh\eta\); choice among (A)--(C) open\\
Cor.~\ref{cor:gue-frame}, v2.29.1 note & the earlier value \(1.197\) was \(\cosh\eta\), the matter--photon Lorentz factor, not \(\gBI\) (r01.1)\\
Prop.~\ref{prop:hawking-gBI-tachyon} (ii) & reciprocal factors put the tachyon on the far side of matter; does not match the axis order (r01)\\
Thm.~\ref{thm:tribonacci} & equivalent to \(\cosh2\eta=e^{\eta}\): matter--tachyon Lorentz factor equals matter--photon Doppler factor (r01.3)\\
\bottomrule
\end{tabular}
\end{center}

\subsection*{v4.9.0-r01 ledger and guard update}
\label{subsec:v490-r01-ledger-guard}

\[
R_{\rm conv}^{v490\text{-}r01}=(0,0,1)
\]
for (provenance of the factor recorded, consistency of the conversion law over three frames decided and its relation to
the tribonacci fixed point established, choice among (A)--(C) open).

The layer spends no guard:
\[
R_{\rm glob}^{(v490\text{-}r01,+)}=(0,1,1,1,1),
\qquad
\Gamma_{\rm guard}^{v490\text{-}r01}=(1,1,1,1,1),
\qquad
\beta_4=1/256\ \text{unspent}.
\]
It does not imply exact finite \(C5\) zero closure, exact all-readout terminality,
\(C6/T2\), a global photon-clock/RH theorem, or global all-readout closure.

% ============================================================================
\clearpage
\section*{v4.9.0-r02 Source-Only: The Round Trip about the Photon and the Umkehrwalze}
\addcontentsline{toc}{section}{v4.9.0-r02 Source-Only: The Round Trip about the Photon and the Umkehrwalze}
\label{sec:v490-r02-round-trip-umkehrwalze}
% ============================================================================

\paragraph{Purpose.}
Theoremlet r19.2 records a \(\mathbb Z_2\) holonomy of the duality circle around the photon.  This layer follows a round
trip matter \(\to\) photon \(\to\) tachyon \(\to\) photon \(\to\) matter in Reading~B, asks on which side of the photon each
passage runs, and asks what the Umkehrwalze of the main body is in physical terms.  It uses the reflector and leakage
language of the OP5 layer (Definition~\ref{def:working-shared-matrix-leakage}).  The computations are
\texttt{llm/k38\_umkehrwalze.py}.  No value is compared with data.  Earlier layers are not rewritten.

\noindent\textbf{Source-only status.}
This v4.9.0-r02 layer is source-only.  It does not build a PDF, does not assign a new DOI,
and does not authorize a public release.  The v4.9.0 Version DOI
\href{https://doi.org/10.5281/zenodo.23534954}{10.5281/zenodo.23534954} and the Concept DOI
\href{https://doi.org/10.5281/zenodo.19210728}{10.5281/zenodo.19210728} are unchanged by it.

\paragraph{Author's motivation (2026-10-21).}
The Umkehrwalze was introduced for a path that does not end at the tachyon but returns, matter \(\to\) photon \(\to\)
tachyon \(\to\) photon \(\to\) matter, read as a parallel transport over the sphere.

\subsection*{v4.9.0-r02 the two sides of the photon}
\label{subsec:v490-r02-sides}

Throughout, \(z=E+iB\), \(s=z\cdot z=E^2-B^2+2i\,E\cdot B\), matter \(s>0\), tachyon \(s<0\), photon \(s=0\)
(\S\ref{subsec:v480-r19-transition-group}).

\paragraph{Theoremlet v4.9.0-r02.1.}
A passage from matter to tachyon at fixed \(|s|\) runs through the upper or the lower half of the \(s\)-plane and carries
\(z\mapsto iz\) or \(z\mapsto -iz\).  The side is the sign of \(\operatorname{Im}s=2\,E\cdot B\); it is Lorentz-invariant,
odd under \(P\) and \(T\), and even under \(C\).  For the round trip:
\begin{enumerate}[leftmargin=2em]
\item[(a)] out and back on the same side (the return retraces the outward path): \(z\mapsto z\);
\item[(b)] out on one side and back on the other (the duality angle keeps its sense, \(\pi/2+\pi/2\)): \(z\mapsto -z\).
\end{enumerate}
On the spinor of Theoremlet r06.4 the round trip (b) is \(\psi\mapsto\pm i\psi\); two round trips give
\(\psi\mapsto-\psi\) with \(z\) restored.

\paragraph{Proof sketch.}
\(s(e^{i\alpha}z)=e^{2i\alpha}s\), so \(\alpha\in[0,\pi/2]\) moves \(s\) through the upper half-plane and
\(\alpha\in[-\pi/2,0]\) through the lower one.  \(SO(3,\mathbb C)\) preserves \(s\) and hence \(\operatorname{Im}s\).  By the
parities of \(E\) and \(B\) (\(P\): \(E\) odd, \(B\) even; \(T\): \(E\) even, \(B\) odd; \(C\): both odd), \(E\cdot B\) is
\(P\)- and \(T\)-odd and \(C\)-even.  Case (b) is a loop around \(s=0\), case (a) is not (Theoremlet r19.2).  The spinor
statement is \(z(e^{i\varphi}\psi)=e^{2i\varphi}z(\psi)\).  Checked numerically.  \(\square\)

\paragraph{Reading.}
The ``double tunnel point'' of the round trip is one point with two sides: the photon, passed on the side
\(E\cdot B>0\) or \(E\cdot B<0\).  The side is the sign of the density of the electromagnetic \(\theta\)-term
\(F\tilde F\).  Only a round trip that changes side encloses the photon.

\subsection*{v4.9.0-r02 the round-trip Umkehrwalze is charge conjugation}
\label{subsec:v490-r02-charge-conjugation}

\paragraph{Theoremlet v4.9.0-r02.2.}
On the field strength the holonomy \(z\mapsto-z\) of the round trip (b) is charge conjugation \(C\), equals \(PT\)
(with \(x\mapsto-x\)), and equals the duality rotation by \(\pi\).  It commutes with \(SO(3,\mathbb C)\) and with the duality
circle, and fixes no non-zero field.

\paragraph{Proof sketch.}
\(C\): \(E\mapsto-E\), \(B\mapsto-B\).  \(T\): \(z\mapsto\bar z\); \(P\): \(z\mapsto-\bar z\); composed, \(z\mapsto-z\).
\(e^{i\pi}=-1\) is central.  \(\square\)

\paragraph{Consequence.}
The Umkehrwalze of the round trip is not an added element.  It is forced by the two-to-one relation between \(z\) and
\(s\), whose branch point is the photon, and in physical terms it is \(C\).  Being central, it cannot serve as the \(J\) of
the readout projector \(C_{\HC}=\tfrac12(\mathrm{id}+J)\) of Definition~\ref{def:working-shared-matrix-leakage}: for
\(J=-1\) that projector is zero.

\subsection*{v4.9.0-r02 the involutive Umkehrwalze must contain time reversal}
\label{subsec:v490-r02-antilinear}

The main body uses \(J\) as an involution that exchanges matter and tachyon (Theorem~\ref{thm:core}, where the inside
view runs with time reversed).  In Reading~B the quarter turn \(z\mapsto iz\) does this exchange, but its square is
\(-1\), so it is not an involution.

\paragraph{Theoremlet v4.9.0-r02.3.}
\begin{enumerate}[leftmargin=2em]
\item[(a)] No \(\mathbb C\)-linear involution \(L\) of \(\mathbb C^3\) satisfies \(s(Lz)=-s(z)\).
\item[(b)] The \(\mathbb C\)-antilinear involutions with \(s(Lz)=-\overline{s(z)}\) are exactly \(Lz=iR\bar z\) with
\(R\in O(3,\mathbb C)\) and \(R\bar R=1\).  Instances are \(R=1\), real orthogonal involutions, and squares of boosts.
\item[(c)] The instance \(R=1\),
\[
J_B z:=i\bar z,\qquad\text{that is}\quad E\mapsto B,\ B\mapsto E,
\]
is the duality rotation by \(\pi/2\) after time reversal.  It exchanges matter and tachyon, maps the null cone to
itself, and preserves the side \(\operatorname{sign}\operatorname{Im}s\).
\end{enumerate}

\paragraph{Proof sketch.}
(a) \(M:=-iL\) preserves \(s\), so \(M\in O(3,\mathbb C)\); \(L^2=1\) gives \(M^2=-1\) and \(\det(M)^2=-1\), impossible.
(b) \(Lz=A\bar z\) with \(A^{\mathsf T}A=-1\) gives \(A=iR\), \(R\in O(3,\mathbb C)\); \(L^2=A\bar A=R\bar R\).
(c) \(i(E-iB)=B+iE\), and \(s\mapsto-\bar s\).  Checked numerically.  \(\square\)

\paragraph{Consequence.}
In three complex dimensions an involutive exchange of matter and tachyon cannot be linear.  It has to be antilinear,
and antilinearity is time reversal.  This matches the inside view of Theorem~\ref{thm:core}, which runs with time
reversed.  The physical reading of the involutive Umkehrwalze is therefore electric--magnetic exchange, a duality
quarter turn composed with \(T\).  It is a symmetry of the source-free Maxwell equations.  Unlike \(C\) it is not
Lorentz-invariant: \(J_B\Lambda=\bar\Lambda J_B\), so it commutes with real rotations only and singles out a frame.

\paragraph{Theoremlet v4.9.0-r02.4.}
The eigen-channels of \(J_B\) lie one on each side of the photon:
\[
J_Bz=z\iff E=B,\ s=2i|E|^2;\qquad J_Bz=-z\iff E=-B,\ s=-2i|E|^2.
\]
The projector \(C_{\HC}=\tfrac12(\mathrm{id}+J_B)\) maps every field to \(\operatorname{Re}s=0\),
\(\operatorname{Im}s\ge0\).  Its image meets the null cone only in \(z=0\).

\paragraph{Proof sketch.}
Direct substitution; checked on 1000 random fields (\(\max|\operatorname{Re}s|=0\), \(\min\operatorname{Im}s>0\)).
\(\square\)

\paragraph{Consequence for Theorem~\ref{thm:core}.}
The compensator \(\tfrac12(f+Jf)\) ``produces the Hawking photon at \(\beta=0\)''.  With \(J=J_B\) it produces the value
\(0\) of the matter--tachyon coordinate \(\operatorname{Re}s=E^2-B^2\), not a null field.  The output keeps
\(E\cdot B\neq0\) and lies on one side of the photon, which is the endpoint of its image.  The statement holds for the real
part of \(s\) only.

\paragraph{Theoremlet v4.9.0-r02.5.}
\(J_B\) reverses the sense of the duality circle, \(J_BR_\alpha=R_{-\alpha}J_B\) with \(R_\alpha z=e^{i\alpha}z\).  Hence
\[
R_\alpha^{-1}J_BR_\alpha=R_{-2\alpha}J_B,
\qquad
R_{\pi/2}^{-1}J_BR_{\pi/2}=-J_B,
\qquad
R_{\pi/4}^{-1}J_BR_{\pi/4}=T .
\]
Passing from matter to the tachyon, reflecting there, and passing back gives the direct reflection followed by the
round-trip Umkehrwalze \(C\).

\paragraph{Proof sketch.}
\(e^{-i\alpha}\,i\,\overline{e^{i\alpha}z}=ie^{-2i\alpha}\bar z\).  Checked numerically.  \(\square\)

\paragraph{Reading.}
This is the Enigma pattern of the computational interpretation surface: a reflector between a forward and a backward
passage.  The two Umkehrwalzen of this layer are related by it.  The central one, \(C\), is the holonomy of a passage
pair.  The involutive one, \(J_B\), is the reflector, and conjugating it by the matter--tachyon passage multiplies it by
\(C\).

\subsection*{v4.9.0-r02 what the OP5 layer contributes}
\label{subsec:v490-r02-op5}

Inserting \(J=J_B\) into the leakage diagnostics \(E_{\rm comm}=[\mathcal M,J]\) and
\(E_{\rm leak}=(\mathrm{id}-C_{\HC})\mathcal MC_{\HC}\) of Definition~\ref{def:working-shared-matrix-leakage} gives the
following, with the transports of Reading~B in place of \(\mathcal M\) (relative norms on a random field):
\begin{center}
\small
\begin{tabular}{@{}lrr@{}}
\toprule
transport \(\mathcal M\) & \(|E_{\rm comm}z|/|z|\) & \(|E_{\rm leak}z|/|C_{\HC}z|\)\\
\midrule
real rotation & \(0\) & \(0\)\\
boost, \(\eta=0.6174\) & \(0.985\) & \(0.526\)\\
duality rotation, \(\alpha=0.4\) & \(0.779\) & \(0.389\)\\
\bottomrule
\end{tabular}
\end{center}
By Proposition~\ref{prop:working-stable-channel-criterion} the admissible channel is stable under the compact part of
the transport only.  Boosts leak because \(J_B\) singles out a frame.  Duality rotations leak because \(J_B\) reverses
their sense (Theoremlet r02.5).  Boosts preserve the side of the field as a whole, so this leakage is into the
\(J_B\)-odd component and does not change side.

\paragraph{Correspondence, offered and not adopted.}
The local OP5 closure theorem works under the identification of the OP5 near-kernel branch with the
\emph{photon-frame HC-cone/reflector route}.  Its halo invariance off tunnel candidates asks for a turning-point
readout before a through-collapse exit.  In the language of this layer a turning point is case (a) of Theoremlet r02.1,
a return on the same side with trivial holonomy, and a through-exit is case (b), a change of side with holonomy \(C\).
The reflector of that route would be antilinear by Theoremlet r02.3.  The radial and \(\HC\)-residual channels
\(\rho_r,\rho_h\), the kernel \(K\) and the tunnel selector have no counterpart in Reading~B yet.  The correspondence is
therefore recorded and not used.  It does not touch the guarded scope of the OP5 closure.

\subsection*{v4.9.0-r02 physical summary in plain terms}
\label{subsec:v490-r02-plain}

The following restates Theoremlets r02.1--r02.5 without notation.  It relies on the working assignment of r07
(matter electric-dominated, tachyon magnetic-dominated), which is a model decision, and on Reading~B, which is
kinematics without dynamics (r10).
\begin{enumerate}[leftmargin=2em]
\item \emph{Two ways past the photon.}  A field changes from matter type to tachyon type only by passing the photon.  It
can pass on one of two sides, according to whether its electric and magnetic parts lean towards each other or away from
each other (the sign of \(E\cdot B\)).  Every observer agrees on the side.
\item \emph{The round trip.}  Going out and coming back on the same side changes nothing.  Going out on one side and
coming back on the other takes the field once around the photon, and it returns with both of its parts reversed.  That
is the field of opposite charges: the round trip acts as charge conjugation.  Only a second round trip restores the
underlying spinor, as with a M\"obius strip.
\item \emph{The mirror between matter and tachyon.}  A map that exchanges matter and tachyon directly and gives the
identity when applied twice exists only if it also reverses time.  The simplest such map exchanges the electric and the
magnetic field.  For light in vacuum this is a known symmetry, electric--magnetic duality combined with time reversal.
The time-reversed inside view of Theorem~\ref{thm:core} has this structure.
\item \emph{What the compensator produces.}  Averaging a field with its mirror image makes the electric and the
magnetic part equally strong, which is the line between matter and tachyon.  The two parts then point the same way.
For light they are perpendicular.  The average is therefore balanced, but it is not a light field.
\item \emph{The two mirror channels.}  The fields that the mirror leaves unchanged lie on one side of the photon, and
the fields it reverses lie on the other.  The two sides of item~1 are the two channels of the mirror.
\end{enumerate}
No statement in this list is compared with data.

\subsection*{v4.9.0-r02 register entries}
\label{subsec:v490-r02-register}

\begin{center}
\small
\begin{tabular}{@{}>{\raggedright\arraybackslash}p{0.34\textwidth}>{\raggedright\arraybackslash}p{0.58\textwidth}@{}}
\toprule
statement & status after v4.9.0-r02\\
\midrule
Umkehrwalze of the round trip & holonomy of a side-changing passage pair, equal to \(C=PT\) on field strengths (r02.1, r02.2)\\
involutive \(J\) exchanging matter and tachyon & must be antilinear (contains \(T\)); representative \(J_B\): \(E\leftrightarrow B\) (r02.3); frame-dependent\\
Thm.~\ref{thm:core}, photon from \(\tfrac12(f+Jf)\) & with \(J_B\): \(\operatorname{Re}s=0\), not a null field; holds for the real part of \(s\) only (r02.4)\\
Def.~\ref{def:working-shared-matrix-leakage} with \(J_B\) & \(E_{\rm comm}=E_{\rm leak}=0\) for real rotations only (r02)\\
OP5 turning point / through-exit & correspondence to cases (a)/(b) of r02.1 recorded, not adopted\\
\bottomrule
\end{tabular}
\end{center}

\subsection*{v4.9.0-r02 ledger and guard update}
\label{subsec:v490-r02-ledger-guard}

\[
R_{\rm uw}^{v490\text{-}r02}=(0,0,1)
\]
for (round-trip Umkehrwalze identified with \(C\), involutive Umkehrwalze characterised as antilinear,
identification with the OP5 halo/tunnel structure open).

The layer spends no guard:
\[
R_{\rm glob}^{(v490\text{-}r02,+)}=(0,1,1,1,1),
\qquad
\Gamma_{\rm guard}^{v490\text{-}r02}=(1,1,1,1,1),
\qquad
\beta_4=1/256\ \text{unspent}.
\]
It does not imply exact finite \(C5\) zero closure, exact all-readout terminality,
\(C6/T2\), a global photon-clock/RH theorem, or global all-readout closure.

% ============================================================================
\clearpage
\section*{v4.9.0-r03 Source-Only: The Compensator as a Two-Arm Interferometer}
\addcontentsline{toc}{section}{v4.9.0-r03 Source-Only: The Compensator as a Two-Arm Interferometer}
\label{sec:v490-r03-compensator-interferometer}
% ============================================================================

\paragraph{Purpose.}
Theoremlet r02.4 found that the compensator \(C(f)=\tfrac12(f+Jf)\) with the involution \(J_Bz=i\bar z\) of r02.3 does
not produce a null field.  This layer looks at the compensator itself in Reading~B: as the output of an interferometer,
on a static charge, on a light wave, and in its product form.  The computations are \texttt{llm/k39\_kompensator.py}.
No value is compared with data.  Earlier layers are not rewritten.

\noindent\textbf{Source-only status.}
This v4.9.0-r03 layer is source-only.  It does not build a PDF, does not assign a new DOI,
and does not authorize a public release.  The v4.9.0 Version DOI
\href{https://doi.org/10.5281/zenodo.23534954}{10.5281/zenodo.23534954} and the Concept DOI
\href{https://doi.org/10.5281/zenodo.19210728}{10.5281/zenodo.19210728} are unchanged by it.

\paragraph{Author's motivation (2026-10-21).}
The picture behind the compensator was the Michelson--Morley interferometer with its half-silvered mirror; a mirror
for the tachyon would act as a periscope that lays both images over each other.

\subsection*{v4.9.0-r03 the two output ports}
\label{subsec:v490-r03-ports}

\paragraph{Theoremlet v4.9.0-r03.1.}
Split a field \(f\) at a beam splitter.  Send one half through the identity and one through \(J\), then recombine.  The
two output ports carry \(C(f)=\tfrac12(f+Jf)\) and \(A(f)=\tfrac12(f-Jf)\), the compensator and its antisymmetric
partner of Definition~r04.1.  The energy splits without loss, \(|Cf|^2+|Af|^2=|f|^2\).  For antilinear \(J\) (r02.3) both
maps are real-linear only, and
\[
A(f)=-i\,C(if),\qquad \operatorname{ran}A=i\,\operatorname{ran}C .
\]
The two ports differ by a duality quarter turn.

\paragraph{Proof sketch.}
\(C+A=1\), \(C^2=C\), \(CA=0\) (Theoremlet r04.2), and \(\operatorname{Re}\langle Cf,Af\rangle=0\) for anti-unitary \(J\).
\(C(if)=\tfrac12(if-iJf)=iA(f)\).  Checked numerically.  \(\square\)

\paragraph{Consequence.}
The bosonic/fermionic split \(C(f)\in\Hplusplus\), \(A(f)\in\Hplusminus\) of the Pauli item in
\S\ref{sec:central} becomes, with an antilinear \(J\), a split into two real subspaces that are images of each other
under multiplication by \(i\).  That is a split into real and imaginary parts, not two independent sectors.  Remark A2 of
r04 already noted that \(\tfrac12(1\pm J)\) is only real-linear for antilinear \(J\).

\subsection*{v4.9.0-r03 a charge and a light wave in the interferometer}
\label{subsec:v490-r03-charge-light}

\paragraph{Theoremlet v4.9.0-r03.2.}
\begin{enumerate}[leftmargin=2em]
\item[(a)] For the Coulomb field of a charge \(q\) (matter type, \(s>0\)), \(J_Bf\) is the field of a magnetic charge
\(q\) (tachyon type).  \(C(f)\) and \(A(f)\) are the fields of dyons with electric charge \(q/2\) and magnetic charge
\(\pm q/2\), with \(s=\pm\tfrac i2|E_0|^2\).
\item[(b)] For a plane light wave, \(J_B\) reverses the direction of propagation and rotates the polarisation by
\(90^\circ\).  This is what a phase-conjugating mirror (time reversal, r02.3) followed by a polarisation rotator (duality
quarter turn) does.  \(C\) of the wave has \(E\parallel B\) and is not a light wave.
\end{enumerate}

\paragraph{Proof sketch.}
\(J_B\): \(E\mapsto B\), \(B\mapsto E\), and \(C\), \(A\) give \(E=\pm B=\tfrac12E_0\).  For \(E=\hat x\), \(B=\hat y\):
\(J_B\) gives \(E=\hat y\), \(B=\hat x\), \(E\times B=-\hat z\).  Checked numerically.  \(\square\)

\paragraph{Reading.}
For light the mirror of r02.3 is an existing optical device, but on light it does not change the type: \(s=0\) stays
\(0\).  A mirror that turns a matter-type field into a tachyon-type field of the same shape must turn electric charge into
magnetic charge.  Under the assignment of r07 the missing tachyon mirror is therefore, in Reading~B, the missing magnetic
charge.  The periscope image, that is the overlay of a field with its mirror image, is a dyon.  This is the content of the
``photon at \(\beta=0\)'' of Theorem~\ref{thm:hawking-triolic} and Theorem~\ref{thm:core} when \(f\) is a static matter
field.

\subsection*{v4.9.0-r03 what the interferometer reads out}
\label{subsec:v490-r03-readout}

\paragraph{Theoremlet v4.9.0-r03.3.}
Put a duality phase \(\varphi\) into the reference arm, \(f\mapsto e^{i\varphi}f\).  The difference of the two port
energies is
\[
D(\varphi)=\bigl|\tfrac12(e^{i\varphi}f+J_Bf)\bigr|^2-\bigl|\tfrac12(e^{i\varphi}f-J_Bf)\bigr|^2
=\operatorname{Im}\bigl(e^{i\varphi}s\bigr).
\]
The fringe has amplitude \(|s|\) and phase \(\arg s\).  At \(\varphi=0\) it reads \(2E\cdot B\), the side of the photon
(r02.1).  Matter and tachyon give fringes half a period apart, and light gives no fringe.  The visibility is
\[
V=\frac{|s|}{|f|^2}=\sqrt{1-\Bigl(\frac{|E\times B|}{u}\Bigr)^2}=\frac1{\cosh2\eta_f},
\qquad u=\tfrac12(E^2+B^2),
\]
where \(\eta_f\) is the rapidity of the frame in which \(E\parallel B\).  That frame exists for every non-null field.

\paragraph{Proof sketch.}
\(D(\varphi)=\operatorname{Re}\langle e^{i\varphi}f,J_Bf\rangle\) and \(\langle z,i\bar z\rangle=i\bar s\).  Further
\(u^2-|E\times B|^2=\tfrac14|s|^2\).  The frame with \(E\parallel B\) moves along \(E\times B\) with speed \(v\),
\(v/(1+v^2)=|E\times B|/(2u)\), so \(|E\times B|/u=\tanh2\eta_f\).  Checked numerically, including the boost to
\(E\parallel B\).  \(\square\)

\paragraph{Reading.}
In this interferometer light gives a null result, and the visibility of any other field is the inverse Lorentz factor of
its energy transport, \(|E\times B|/u=\tanh2\eta_f\).  The double rapidity \(2\eta_f\) appears unforced: the field's rest
frame moves at \(\tanh\eta_f\), its energy at \(\tanh2\eta_f\).  If \(\eta_f\) were \(\eta_{\rm sym}\), the visibility would be
\(1/\gL=0.536\) and the energy speed \(0.844\).  This is arithmetic.  No field is identified with \(\eta_{\rm sym}\), and an
identification would be a prediction requiring a pre-registration note.

\subsection*{v4.9.0-r03 sum and product}
\label{subsec:v490-r03-product}

Proposition~\ref{prop:compensator} calls \(C(f)\) the projection onto \(K_3=e_1e_2\) and states that it is orthogonal to
both inputs.  A projection is not orthogonal to its input: \(\operatorname{Re}\langle f,Cf\rangle=|Cf|^2\).

\paragraph{Theoremlet v4.9.0-r03.4.}
The product of a field with its mirror image has the properties claimed for \(e_3=e_1e_2\):
\[
z\times J_Bz=2\,E\times B,\qquad z\cdot J_Bz=2iu .
\]
The vector part is twice the Poynting vector, real and orthogonal to \(E\) and to \(B\).  The scalar part is the energy
density.  The pair \((u,E\times B)\) is null, \(|E\times B|=u\), exactly for light.

\paragraph{Proof sketch.}
\(z\times\bar z=-2iE\times B\) and \(z\cdot\bar z=E^2+B^2\).  Checked numerically.  \(\square\)

\paragraph{Consequence.}
The orthogonality in Proposition~\ref{prop:compensator} holds for the product of \(f\) with \(Jf\), not for their
average.  In Reading~B the product yields the energy and momentum of the field, and it is light-like exactly on the
photon sector.  The average yields a dyon.  Which of the two the main body means by the compensator is not decided here.
Both readings are recorded.

\subsection*{v4.9.0-r03 register entries}
\label{subsec:v490-r03-register}

\begin{center}
\small
\begin{tabular}{@{}>{\raggedright\arraybackslash}p{0.34\textwidth}>{\raggedright\arraybackslash}p{0.58\textwidth}@{}}
\toprule
statement & status after v4.9.0-r03\\
\midrule
compensator \(C\), partner \(A\) & the two output ports of a two-arm interferometer with \(J\) in one arm; real-linear, \(A(f)=-iC(if)\) (r03.1)\\
Pauli item, \(\Hplusplus\)/\(\Hplusminus\) split & with antilinear \(J\) a real/imaginary split, the sectors are \(i\)-images of each other (r03.1)\\
Thm.~\ref{thm:hawking-triolic}, Thm.~\ref{thm:core}, photon at \(\beta=0\) & for a static matter field \(C(f)\) is a dyon field (r03.2)\\
Prop.~\ref{prop:compensator}, orthogonal to both inputs & false for the average; true for the product \(z\times J_Bz=2E\times B\) (r03.4)\\
tachyon mirror & on light: phase conjugation plus polarisation rotation; on matter-type fields it needs magnetic charge (r03.2)\\
\bottomrule
\end{tabular}
\end{center}

\subsection*{v4.9.0-r03 ledger and guard update}
\label{subsec:v490-r03-ledger-guard}

\[
R_{\rm comp}^{v490\text{-}r03}=(0,0,1)
\]
for (interferometer form of the compensator established, readout \(D(\varphi)=\operatorname{Im}(e^{i\varphi}s)\) and
visibility \(1/\cosh2\eta_f\) derived, choice between average and product reading open).

The layer spends no guard:
\[
R_{\rm glob}^{(v490\text{-}r03,+)}=(0,1,1,1,1),
\qquad
\Gamma_{\rm guard}^{v490\text{-}r03}=(1,1,1,1,1),
\qquad
\beta_4=1/256\ \text{unspent}.
\]
It does not imply exact finite \(C5\) zero closure, exact all-readout terminality,
\(C6/T2\), a global photon-clock/RH theorem, or global all-readout closure.

% ============================================================================
\clearpage
\section*{v4.9.0-r04 Source-Only: The Compensator as a Matrix on the Electric--Magnetic Doublet}
\addcontentsline{toc}{section}{v4.9.0-r04 Source-Only: The Compensator as a Matrix on the Electric--Magnetic Doublet}
\label{sec:v490-r04-doublet-matrix}
% ============================================================================

\paragraph{Purpose.}
On \(\mathbb C^3\) the involution \(J_B\) of r02.3 is antilinear.  It has no complex matrix, and the compensator built from
it is only real-linear (r03.1).  This layer writes the field as the real pair \((E,B)\) instead.  Then \(J_B\), the
compensator and the readout of r03 become ordinary \(2\times2\) matrices, and the state of a field becomes a
\(2\times2\) matrix of the kind used for polarisation in optics.  The computations are
\texttt{llm/k40\_dublett\_matrix.py}.  No value is compared with data.  Earlier layers are not rewritten.

\noindent\textbf{Source-only status.}
This v4.9.0-r04 layer is source-only.  It does not build a PDF, does not assign a new DOI,
and does not authorize a public release.  The v4.9.0 Version DOI
\href{https://doi.org/10.5281/zenodo.23534954}{10.5281/zenodo.23534954} and the Concept DOI
\href{https://doi.org/10.5281/zenodo.19210728}{10.5281/zenodo.19210728} are unchanged by it.

\paragraph{Author's question (2026-10-21).}
Would it be better to put the compensator into a matrix?

\subsection*{v4.9.0-r04 the doublet and its matrices}
\label{subsec:v490-r04-doublet}

Write \(F=[\,E\mid B\,]\in\mathbb R^{3\times2}\) for \(z=E+iB\), and let
\(\sigma_x=\bigl(\begin{smallmatrix}0&1\\1&0\end{smallmatrix}\bigr)\),
\(\sigma_z=\bigl(\begin{smallmatrix}1&0\\0&-1\end{smallmatrix}\bigr)\), \(R_\alpha\) the rotation of \(\mathbb R^2\) by
\(\alpha\).

\paragraph{Theoremlet v4.9.0-r04.1.}
The operations of r02 act on \(F\) from the right by real \(2\times2\) matrices:
\[
e^{i\alpha}z\leftrightarrow FR_\alpha^{\mathsf T},\qquad
J_Bz=i\bar z\leftrightarrow F\sigma_x,\qquad
T:\ \bar z\leftrightarrow F\sigma_z,\qquad
C:\ -z\leftrightarrow -F .
\]
They generate \(O(2)\): \(\sigma_x\) and \(\sigma_z\) are reflections and \(R_\alpha\) rotations.  The relation
\(J_BR_\alpha=R_{-\alpha}J_B\) of r02.5 is the relation between a reflection and a rotation in \(O(2)\), and the Enigma
composite of r02.5 is the matrix identity \(R_{\pi/2}\sigma_xR_{\pi/2}^{-1}=-\sigma_x\).  Spatial rotations act from the
left, \(F\mapsto RF\), and commute with all of these.  Boosts mix \(E\) and \(B\) through cross products, are not of the
form \(F\mapsto LFM\), and do not commute with \(\sigma_x\).  The compensator and its partner are ordinary orthogonal
projections,
\[
C_{\HC}=\tfrac12(1+\sigma_x)=\tfrac12\begin{pmatrix}1&1\\1&1\end{pmatrix},\qquad
A=\tfrac12(1-\sigma_x)=\tfrac12\begin{pmatrix}1&-1\\-1&1\end{pmatrix},
\]
and \(FC_{\HC}\) is the compensator of r02.4 and r03.

\paragraph{Proof sketch.}
\(e^{i\alpha}(E+iB)=(E\cos\alpha-B\sin\alpha)+i(E\sin\alpha+B\cos\alpha)\); \(i(E-iB)=B+iE\); \(E-iB\).  The boost
formula \(E'=\gamma(E+v\times B)-\dots\) couples the columns through \(v\times\).  Checked numerically.  \(\square\)

\paragraph{Consequence.}
The antilinearity of r02.3 is not lost.  It is the statement that \(\sigma_x\) is a reflection, \(\det\sigma_x=-1\),
while the duality circle consists of rotations.  Over \(\mathbb R\) the compensator is linear, and the obstruction of
r03.1 disappears.

\subsection*{v4.9.0-r04 the state matrix}
\label{subsec:v490-r04-state}

\paragraph{Theoremlet v4.9.0-r04.2.}
The \(2\times2\) Gram matrix \(G=F^{\mathsf T}F\) of \(E\) and \(B\) decomposes as
\[
G=u\,\mathbf 1+\tfrac12\bigl(\operatorname{Re}s\;\sigma_z+\operatorname{Im}s\;\sigma_x\bigr),
\qquad
\det G=|E\times B|^2,\qquad
u^2-\det G=\tfrac14|s|^2 .
\]
Put \(\rho=G/\operatorname{tr}G\) and \(b=(\operatorname{Re}s,\operatorname{Im}s)/(2u)\), so that
\(\rho=\tfrac12(\mathbf 1+b_1\sigma_z+b_2\sigma_x)\).  Then:
\begin{enumerate}[leftmargin=2em]
\item[(a)] \(|b|=1/\cosh2\eta_f\le1\), with \(\eta_f\) as in r03.3.  \(|b|=1\) exactly when \(E\parallel B\), and \(b=0\)
exactly for a null field.  Matter (\(b_1>0\)) and tachyon (\(b_1<0\)) lie on opposite sides of the unit disk, light at its
centre, and the side of the photon of r02.1 is the sign of \(b_2\).
\item[(b)] The duality rotation by \(\alpha\) rotates \(b\) by \(2\alpha\).  \(J_B\) reflects \(b_1\mapsto-b_1\), and
\(T\) reflects \(b_2\mapsto-b_2\).
\item[(c)] Under a boost the traceless part of \(G\), that is \(s\), is unchanged and \(u\) changes, so \(b\) moves along its
own direction.  It reaches \(|b|=1\) in the frame with \(E\parallel B\) and tends to \(0\) as the boost grows.
\end{enumerate}
\(G\) is symmetric and has no \(\sigma_y\) part, so the states fill a disk, not a sphere.

\paragraph{Proof sketch.}
\(G_{11}=E^2\), \(G_{22}=B^2\), \(G_{12}=E\cdot B\); \(\det G=E^2B^2-(E\cdot B)^2\).  (a) follows from r03.3.  (b): \(G\mapsto
R_\alpha GR_\alpha^{\mathsf T}\), and the traceless part of a symmetric matrix turns by twice the angle.  (c): \(s\) is
Lorentz-invariant (r06.3).  Checked numerically, including a boost scan through the frame \(E\parallel B\).  \(\square\)

\paragraph{Reading.}
This is the formalism of Stokes parameters for the doublet \((E,B)\) in place of two polarisation components.  The doubled
angle in (b) is the half-angle of the spinor of r06.4 seen from the other side: a duality rotation by \(\alpha\) moves the
state by \(2\alpha\).  Light is the unpolarised centre, and matter and tachyon are the two fully polarised ends of one
axis.  Boosting any field moves it towards light.

\subsection*{v4.9.0-r04 the interferometer as a measurement}
\label{subsec:v490-r04-measurement}

\paragraph{Theoremlet v4.9.0-r04.3.}
The energy in the compensator port of r03.1 is \(|FC_{\HC}|^2=\operatorname{tr}(GC_{\HC})=u+E\cdot B\).  With the duality phase
\(\varphi\) in the reference arm, the port difference of r03.3 is
\[
D(\varphi)=\operatorname{Im}(e^{i\varphi}s)=2u\;b\cdot(\sin\varphi,\cos\varphi).
\]
The interferometer measures the component of \(b\) along a direction set by \(\varphi\).  The compensator is the
measurement along \(\sigma_x\), and its partner is the opposite outcome.

\paragraph{Proof sketch.}
\(C_{\HC}\) is a symmetric projection, so \(|FC_{\HC}|^2=\operatorname{tr}(C_{\HC}GC_{\HC})=\operatorname{tr}(GC_{\HC})\).
Expand \(\operatorname{Im}(e^{i\varphi}s)\).  Checked numerically.  \(\square\)

\paragraph{Both readings of r03 in one matrix.}
The average reading of r03.4 is the projection \(F\mapsto FC_{\HC}\), linear in \(F\).  The product reading is
quadratic.  Its scalar part is \(\operatorname{tr}G=2u\), its magnitude is \(\sqrt{\det G}=|E\times B|\), and its direction is
the cross product of the two columns of \(F\).  The trace, traceless part and determinant of \(G\) are energy, the
invariant \(s\), and the magnitude of the energy flux.  The \(3\times2\) matrix \(F\) carries in addition the direction of
the flux, which \(G\) does not see.

\subsection*{v4.9.0-r04 relation to the compensator matrix of the main body}
\label{subsec:v490-r04-hc-matrix}

Section~\ref{app:hc-tools} defines a \emph{global Heisenberg-compensator matrix} only as the extension of
\(C=\tfrac12(\mathrm{id}+J)\) to the multi-layer architecture, and a frame-balancing tool that splits \(X=X^++X^-\) by
\(J\).  The doublet supplies concrete instances in Reading~B: \(X^+=FC_{\HC}\) and \(X^-=FA\).  In matrix form the leakage
diagnostics of Definition~\ref{def:working-shared-matrix-leakage}, with \(J=\sigma_x\) acting from the right, read
(relative norms on a random field):
\begin{center}
\small
\begin{tabular}{@{}lrr@{}}
\toprule
transport \(\mathcal M\) & \(|E_{\rm comm}|/|F|\) & \(|E_{\rm leak}|/|FC_{\HC}|\)\\
\midrule
spatial rotation (left) & \(0\) & \(0\)\\
boost, \(\eta=0.6174\) & \(1.113\) & \(0.505\)\\
duality rotation, \(\alpha=0.4\) (right) & \(0.779\) & \(0.389\)\\
\bottomrule
\end{tabular}
\end{center}
This repeats r02 in matrix form.  Left and right multiplication commute, so rotations cannot leak.  Duality rotations
leak because \(R_\alpha\) does not commute with \(\sigma_x\), and boosts leak because they do not factorise.  The values
for the boost differ from the r02 table because the random field differs; vanishing and non-vanishing agree.  Identifying
the Millennium matrix of the main body with a matrix on the doublet is not done here.

\paragraph{Open: the third axis.}
\(G\) is real and symmetric, so the state space is a disk.  A \(\sigma_y\) part would need an antisymmetric component of
the state matrix, for instance from complex amplitudes of time-harmonic fields.  Whether such a third axis, which would
turn the disk into a sphere, has a meaning in the source is not decided here.

\subsection*{v4.9.0-r04 register entries}
\label{subsec:v490-r04-register}

\begin{center}
\small
\begin{tabular}{@{}>{\raggedright\arraybackslash}p{0.34\textwidth}>{\raggedright\arraybackslash}p{0.58\textwidth}@{}}
\toprule
statement & status after v4.9.0-r04\\
\midrule
compensator \(C=\tfrac12(1+J)\) & on the doublet \((E,B)\) the linear projection \(\tfrac12(1+\sigma_x)\); \(J_B\), \(T\), duality, \(C\) generate \(O(2)\) (r04.1)\\
state of a field & Gram matrix \(G=u\mathbf 1+\tfrac12(\operatorname{Re}s\,\sigma_z+\operatorname{Im}s\,\sigma_x)\), \(\det G=|E\times B|^2\); point \(b\) in the unit disk, light at the centre, \(|b|=1/\cosh2\eta_f\) (r04.2)\\
interferometer of r03 & measurement of \(b\) along a direction set by \(\varphi\) (r04.3)\\
global HC matrix (\S\ref{app:hc-tools}) & doublet gives concrete \(X^\pm\); identification with the Millennium matrix not made\\
\bottomrule
\end{tabular}
\end{center}

\subsection*{v4.9.0-r04 ledger and guard update}
\label{subsec:v490-r04-ledger-guard}

\[
R_{\rm mat}^{v490\text{-}r04}=(0,0,1)
\]
for (doublet matrix form of the compensator established, state disk and readout as measurement established,
identification with the global compensator matrix of the main body open).

The layer spends no guard:
\[
R_{\rm glob}^{(v490\text{-}r04,+)}=(0,1,1,1,1),
\qquad
\Gamma_{\rm guard}^{v490\text{-}r04}=(1,1,1,1,1),
\qquad
\beta_4=1/256\ \text{unspent}.
\]
It does not imply exact finite \(C5\) zero closure, exact all-readout terminality,
\(C6/T2\), a global photon-clock/RH theorem, or global all-readout closure.

% ============================================================================
\clearpage
\section*{v4.9.0-r05 Source-Only: The Millennium Matrix on the Doublet, a Third Axis near Rotating Black Holes, and a Photon from Matter and Tachyon}
\addcontentsline{toc}{section}{v4.9.0-r05 Source-Only: The Millennium Matrix on the Doublet, a Third Axis near Rotating Black Holes, and a Photon from Matter and Tachyon}
\label{sec:v490-r05-layer-transport-kerr-photon}
% ============================================================================

\paragraph{Purpose.}
Three questions of the author follow r04.
\begin{enumerate}[leftmargin=2em]
\item Can the Millennium matrix of the main body be the doublet matrix of r04?  The author also asks for a more
meaningful name.
\item Can amplitudes, the missing third axis of r04, enter near a black hole?
\item Can a photon arise when a tachyon meets a forming massive particle, and could such a photon run back and create
the tachyon at its origin?
\end{enumerate}
The computations are \texttt{llm/k41\_schichttransport\_kerr\_photon.py}.  No value is compared with data.  Earlier
layers are not rewritten.

\noindent\textbf{Source-only status.}
This v4.9.0-r05 layer is source-only.  It does not build a PDF, does not assign a new DOI,
and does not authorize a public release.  The v4.9.0 Version DOI
\href{https://doi.org/10.5281/zenodo.23534954}{10.5281/zenodo.23534954} and the Concept DOI
\href{https://doi.org/10.5281/zenodo.19210728}{10.5281/zenodo.19210728} are unchanged by it.

\subsection*{v4.9.0-r05 the Millennium matrix on the doublet}
\label{subsec:v490-r05-millennium}

The main body asks three things of the Millennium matrix \(\mathcal M\).  It transports one object through its layers
(\S\ref{sec:millennium-hc}).  It commutes with the involution and with the compensator, condition (T4) of
Definition~\ref{def:tensor-structural-kernel}.  And the compensator is its spectral projector
(Theorem~\ref{thm:hc-spectral-projector}).

\paragraph{Theoremlet v4.9.0-r05.1.}
On the doublet of r04, with \(J=\sigma_x\) and \(C_{\HC}=\tfrac12(1+\sigma_x)\) acting from the right:
\begin{enumerate}[leftmargin=2em]
\item[(a)] The real \(2\times2\) matrices satisfying (T4) are exactly \(M=\lambda_+C_{\HC}+\lambda_-A\).  Together with
spatial rotations acting from the left, these are all transports of the doublet compatible with (T4).
\item[(b)] For \(\lambda_+\neq\lambda_-\), \(C_{\HC}\) is the spectral projector of \(M\) onto the eigenvalue \(\lambda_+\),
as Theorem~\ref{thm:hc-spectral-projector} requires.
\item[(c)] Of the transports of Reading~B, (T4) holds for spatial rotations, for \(J_B\) itself and for the duality
rotation by \(\pi\), that is \(C\).  It fails for every other duality rotation, for \(T\) and for boosts (r02, r04).
\item[(d)] The invertible members with determinant \(1\) are \(e^{t\sigma_x}=e^{t}C_{\HC}+e^{-t}A\).  They preserve
\(\operatorname{Re}s=E^2-B^2\), the matter--tachyon coordinate, and weight the two ports of r03 by \(e^{\pm t}\).
\end{enumerate}

\paragraph{Proof sketch.}
\(M\sigma_x=\sigma_xM\) is a linear condition with a two-dimensional solution space spanned by \(1\) and \(\sigma_x\).
\(e^{t\sigma_x}{}^{\mathsf T}\sigma_ze^{t\sigma_x}=\sigma_z\) since \(\sigma_x\sigma_z=-\sigma_z\sigma_x\).  Checked
numerically.  \(\square\)

\paragraph{Consequence.}
On the doublet the Millennium matrix can only be a \emph{port-weighting matrix}.  It leaves the compensator port and the
partner port invariant and multiplies them by two factors.  It is not the Lorentz transport of Reading~B, since boosts
violate (T4).  The doublet therefore realises the compensator part of \(\mathcal H_{16}\), not the full \(\mathcal M_{16}\).
The four-term layer decomposition \(\mathcal M_{\rm tri}+\mathcal M_{\rm alg}+\mathcal M_{\rm cl}+\mathcal M_{\rm proj}\)
and the index \(16\) have no counterpart on a two-dimensional doublet.  With \(t=\eta_{\rm sym}\) the port weights would be
the Doppler factors \(e^{\pm\eta}\) of way~(C) in r01.  That is arithmetic, not an identification.

\paragraph{Naming (author's request, 2026-10-21).}
From this layer on, \(\mathcal M_{16}\) is called the \emph{layer-transport matrix}.  The name says what
\S\ref{sec:millennium-hc} defines: one object carried through its layers.  Its restriction to the doublet is called the
\emph{port-weighting matrix}.  Symbols, labels and the wording of earlier layers stay unchanged, so ``Millennium matrix''
in earlier text means the same object.  A renaming of the earlier text and of the PDF keywords would rewrite released
layers.  It is not done here.

\subsection*{v4.9.0-r05 a third axis, and rotating black holes}
\label{subsec:v490-r05-third-axis}

\paragraph{Theoremlet v4.9.0-r05.2.}
\begin{enumerate}[leftmargin=2em]
\item[(a)] At a point, \(E\) and \(B\) are real, so the state matrix of r04 has no \(\sigma_y\) part anywhere, near a black
hole or not.
\item[(b)] For time-harmonic fields with complex amplitudes \(\hat E,\hat B\), the matrix of time averages
\(\hat G=\hat F^\dagger\hat F\) is Hermitian, and its \(\sigma_y\) part is \(-2\operatorname{Im}(\hat E^*\cdot\hat B)\).  For a plane
wave this is \(\pm1\) (relative to \(\operatorname{tr}\hat G\)) for the two circular polarisations and \(0\) for linear
polarisation.  The third axis is helicity, the quantity generated by the duality rotation.  With it the disk of r04
becomes a sphere, with light on the helicity axis.
\item[(c)] The gravitational field has the same structure.  Its Weyl curvature is a complex symmetric \(3\times3\) matrix,
gravito-electric plus \(i\) gravito-magnetic.  For a Schwarzschild black hole the only component
\(\Psi_2=-M/r^3\) is real.  For a Kerr black hole \(\Psi_2=-M/(r-ia\cos\theta)^3\) is complex.  \(\operatorname{Im}\Psi_2^2\), the
Pontryagin part of the curvature, is positive in one hemisphere, zero on the equator and negative in the other.
\end{enumerate}

\paragraph{Proof sketch.}
(b): \(\operatorname{tr}(\hat G\sigma_y)=-2\operatorname{Im}\hat G_{12}\); for \(\hat B=\hat z\times\hat E\) with
\(\hat E=(1,\pm i,0)/\sqrt2\) one gets \(\mp1\) up to orientation.  (c): standard form of \(\Psi_2\) in the Kinnersley tetrad;
the sign of \(\operatorname{Im}\Psi_2^2\) follows the sign of \(\cos\theta\).  Checked numerically for \(M=1\), \(a=0.9\),
\(r=3\).  \(\square\)

\paragraph{Reading.}
Amplitudes do not enter because a black hole is near.  They enter with waves and quantum states.  What a black hole adds
depends on its rotation.  A static one is mirror-symmetric and singles out neither helicity.  A rotating one puts its own
curvature off the real axis, on opposite sides in its two hemispheres, which is the gravitational counterpart of the two
sides of the photon in r02.1.  Whether this couples to the helicity axis of (b), and whether Hawking radiation of a
rotating black hole would then carry a helicity imbalance between the hemispheres, is open.  A quantitative statement is a
prediction and requires a pre-registration note.  Theorem~\ref{thm:core} treats the static case only.

\subsection*{v4.9.0-r05 a photon from matter and tachyon}
\label{subsec:v490-r05-photon}

\paragraph{Author's question (2026-10-21).}
A tachyon meets a massive particle that is just forming.  Instead of reaching a massive state, the encounter yields a
photon.  The tachyon develops a pull through the impact, so the photon is drawn along the tachyon's former path,
collides again and again, reaches the tachyon's origin, vanishes there, and thereby creates the tachyon.

\paragraph{Theoremlet v4.9.0-r05.3.}
\begin{enumerate}[leftmargin=2em]
\item[(a)] Let \(z_M\) be of matter type and \(z_T\) of tachyon type with \(s(z_T)=-s(z_M)\).  Then
\(s(z_M+z_T)=2\,z_M\cdot z_T\).  The superposition is a null field exactly when \(z_M\cdot z_T=0\).  For a static electric
field \(E\) and a static magnetic field \(B\) of equal strength, the superposition is null exactly when \(E\perp B\).  Its
energy then flows at the speed of light along \(E\times B\), perpendicular to both.
\item[(b)] At an angle \(\vartheta\) between \(E\) and \(B\), the superposition sits at \(b=(0,\cos\vartheta)\) in the disk of
r04.  \(\vartheta=90^\circ\) is the photon at the centre, and \(\vartheta=0\) is the dyon of r03.2 at the edge.
\item[(c)] Let a tachyon leave the origin \(O\) and meet the particle at \(K\), a distance \(d\) away, at speed \(v>1\).  A
photon running back from \(K\) to \(O\) arrives \(d(1+1/v)\) after the departure, and this arrival lies in the future light
cone of the departure in every frame.  The order of departure and meeting is frame-dependent (for observer speeds above
\(1/v\) it reverses), but the photon leg is not.
\end{enumerate}

\paragraph{Proof sketch.}
(a) Expand \(s\).  For \(z_M=E\), \(z_T=iB\): \(z_M\cdot z_T=iE\cdot B\).  Null fields have \(|E\times B|=u\) (r03.3).
(b) \(s=2i|E|^2\cos\vartheta\) and \(u=|E|^2\).  (c) The separation of departure \((0,0)\) and photon arrival
\((d/v+d,0)\) is timelike and future-directed; the tachyon leg is spacelike.  Checked numerically for \(v=1/\bsym\).
\(\square\)

\paragraph{Reading.}
The first part of the picture has a precise form.  Matter and tachyon can combine to a photon, but only with crossed
orientation.  With parallel orientation they give the dyon of r03.  The photon does not move along the tachyon's path.
It moves perpendicular to both inputs.  The last part, a photon that arrives at the origin and thereby creates the
tachyon, cannot hold for a photon running forward in time: it always arrives after the tachyon has left.  What does hold
is the familiar reinterpretation of tachyon exchange.  Since the tachyon leg is spacelike, ``emitted at \(O\), absorbed at
\(K\)'' and ``emitted at \(K\), absorbed at \(O\)'' are the same exchange seen by different observers.  A closed
self-creation would need a photon leg running backward in time.  Without one, the repeated encounter of the picture is a
cycle of period \(d(1+1/v)\), not a self-creation.  If such a cycle returns on the other side of the photon, the tachyon
created in the next cycle is the charge conjugate of the first (r02.1).  Reading~B has no dynamics (r10), so the pull of
the picture cannot be derived or excluded here.

\subsection*{v4.9.0-r05 register entries}
\label{subsec:v490-r05-register}

\begin{center}
\small
\begin{tabular}{@{}>{\raggedright\arraybackslash}p{0.34\textwidth}>{\raggedright\arraybackslash}p{0.58\textwidth}@{}}
\toprule
statement & status after v4.9.0-r05\\
\midrule
Millennium matrix, (T4), Thm.~\ref{thm:hc-spectral-projector} & on the doublet only \(\lambda_+C_{\HC}+\lambda_-A\) (port weighting); boosts and general duality rotations violate (T4); full \(\mathcal M_{16}\) not identified (r05.1)\\
name & \(\mathcal M_{16}\): layer-transport matrix; doublet restriction: port-weighting matrix; earlier wording unchanged\\
third axis of r04 & helicity, for complex amplitudes; Kerr curvature off the real axis with opposite sides in the two hemispheres; coupling open (r05.2)\\
photon from matter \(+\) tachyon & null exactly for crossed orientation, flux perpendicular to both; return photon cannot create the tachyon (r05.3)\\
\bottomrule
\end{tabular}
\end{center}

\subsection*{v4.9.0-r05 ledger and guard update}
\label{subsec:v490-r05-ledger-guard}

\[
R_{\rm lt}^{v490\text{-}r05}=(0,0,1)
\]
for (doublet form of the layer-transport matrix decided, third axis and photon question answered structurally,
identification of the full \(\mathcal M_{16}\) and the Kerr coupling open).

The layer spends no guard:
\[
R_{\rm glob}^{(v490\text{-}r05,+)}=(0,1,1,1,1),
\qquad
\Gamma_{\rm guard}^{v490\text{-}r05}=(1,1,1,1,1),
\qquad
\beta_4=1/256\ \text{unspent}.
\]
It does not imply exact finite \(C5\) zero closure, exact all-readout terminality,
\(C6/T2\), a global photon-clock/RH theorem, or global all-readout closure.

% ============================================================================
\clearpage
\section*{v4.9.0-r06 Source-Only: The Tachyon Hit Chain, its Acceleration, and the Higgs Field}
\addcontentsline{toc}{section}{v4.9.0-r06 Source-Only: The Tachyon Hit Chain, its Acceleration, and the Higgs Field}
\label{sec:v490-r06-hit-chain-higgs}
% ============================================================================

\paragraph{Purpose.}
Theoremlet r05.3(c) stated that a photon returning to the origin of a tachyon cannot create it.  The author objected that the
tachyon already runs backward in time, and then refined the picture.  The Higgs field tries to create matter along a line.
A tachyon that has been travelling along this line for a long time hits it again and again.  Each hit is a point of
spacetime at which an essentially identical photon arises.  This layer makes r05.3(c) precise, works out the hit chain,
asks what an accelerating or infinitely accelerated tachyon does to the hits, and checks the picture against the Higgs
field.  Special relativity is used with \(c=1\).  The computations are \texttt{llm/k42\_tachyon\_trefferkette.py}.  No
value is compared with data.  Earlier layers are not rewritten.

\noindent\textbf{Source-only status.}
This v4.9.0-r06 layer is source-only.  It does not build a PDF, does not assign a new DOI,
and does not authorize a public release.  The v4.9.0 Version DOI
\href{https://doi.org/10.5281/zenodo.23534954}{10.5281/zenodo.23534954} and the Concept DOI
\href{https://doi.org/10.5281/zenodo.19210728}{10.5281/zenodo.19210728} are unchanged by it.

\subsection*{v4.9.0-r06 closing the loop, made precise}
\label{subsec:v490-r06-loop}

\paragraph{Theoremlet v4.9.0-r06.1.}
\begin{enumerate}[leftmargin=2em]
\item[(a)] A photon needs the time \(d\) for the distance \(d\).  A tachyon of speed \(v>1\) covers the same distance in the
time \(d/v\), and even read backward it takes back only \(d/v<d\).  One photon leg and one tachyon leg therefore never close
a loop.  The remaining gap is \(d(1-1/v)\), which is \(0.451\,d\) for \(v=1/\bsym\).
\item[(b)] Two tachyon legs emitted in frames moving relative to each other can close or overshoot a loop.  A return leg
that is instantaneous in a frame of speed \(u\) arrives at the origin at \(d/v-ud\), before the start for \(u>1/v\).
\item[(c)] For \(v=1/\bsym\) and \(u=\bsym\) the return arrives exactly at the start.  It runs back along the same
spacelike line, because the outgoing leg is itself instantaneous in that frame.  This is one exchange read in two
directions, case (a) of Theoremlet r02.1, and carries no information around a loop.
\end{enumerate}

\paragraph{Proof sketch.}
Arithmetic on the worldlines; the return leg is the line of simultaneity \(t-t_B=u(x-x_B)\) through
\(B=(d/v,d)\).  Checked numerically.  \(\square\)

\paragraph{Correction of wording in r05.}
The phrase ``a photon leg running backward in time'' in the reading of r05.3 is made precise by (a) and (b).  What fails
is one null leg plus one tachyon leg.  Two tachyon legs do not fail.  The statement r05.3(c) itself stands.

\subsection*{v4.9.0-r06 the hit chain at constant speed}
\label{subsec:v490-r06-chain}

\paragraph{Theoremlet v4.9.0-r06.2.}
Let hits occur along a line at speed \(v>1\), each emitting an identical light signal.  The energy of each signal comes
from the field at the hit, so nothing is carried faster than light.
\begin{enumerate}[leftmargin=2em]
\item[(a)] The signals superpose to a cone, the envelope of their Huygens wavelets, with \(\cos\vartheta=1/v\).  For
\(v=1/\bsym\) this gives \(\cos\vartheta=\bsym\), \(\vartheta=56.68^\circ\).
\item[(b)] The hits are spacelike separated.  Observers slower than \(1/v\) see them run in the direction of the
tachyon, observers faster than \(1/v\) see them run the other way, and the observer of speed exactly \(u=1/v\) sees all hits
at once.  For \(v=1/\bsym\) that observer moves at \(\bsym\), the matter-frame offset.
\item[(c)] The superluminal rapidity of r20.2 is \(\operatorname{artanh}v=\eta+i\pi/2\).  The boost to the observer of (b)
removes \(\eta\) and leaves \(i\pi/2\), the duality quarter turn of r07.2.  In that frame the tachyon is no longer a moving
object but a phase.  The plane wave emitted there perpendicular to the line appears in the original frame at
\(\cos\vartheta'=\bsym\), which is the cone of (a).
\item[(d)] With imaginary mass \(i\mu\), the tachyon at \(v=1/\bsym\) has energy \(\mu\sinh\eta\) and momentum
\(\mu\cosh\eta=1.197\,\mu\), the neighbour cross term of r18.
\end{enumerate}

\paragraph{Proof sketch.}
(a) Minimal arrival time over emission points; checked by a discrete Huygens sum.  (b)
\(\Delta t'=\gamma(\Delta t-u\Delta x)\) with \(\Delta t=\Delta x/v\).  (c) \(\tanh(x+i\pi/2)=\coth x\) and the aberration
formula with \(\vartheta=90^\circ\).  (d) \(E=\mu/\sqrt{v^2-1}\), \(p=\mu v/\sqrt{v^2-1}\).  \(\square\)

\paragraph{Reading.}
``The tachyon was already on its way before the experiment started'' is the frame dependence of (b).  There is an
observer for whom the whole path exists at one instant.  For \(v=1/\bsym\) this observer is the matter frame.  A
superluminal chain of locally triggered emissions is the situation of a light spot sweeping faster than light across a
screen.  It is permitted because each point emits by itself.

\subsection*{v4.9.0-r06 acceleration and infinite speed}
\label{subsec:v490-r06-acceleration}

\paragraph{Theoremlet v4.9.0-r06.3.}
\begin{enumerate}[leftmargin=2em]
\item[(a)] \emph{Constant proper acceleration.}  The tachyonic counterpart of hyperbolic motion is
\(t=R\cosh\sigma\), \(x=R\sinh\sigma\), with proper acceleration \(1/R\) everywhere.  Its speed \(v=\coth\sigma\) and its
coordinate acceleration \(dv/dt\) diverge at \(\sigma=0\), that is \(t=R\), \(x=0\).  The hit rate per unit time diverges
there as well, so the hits near that point arrive at once as a flash.  Beyond \(\sigma=0\) the parameter runs backward in
time.  By reinterpretation that part is a partner moving forward, and the point \((R,0)\) is a pair event.  A signal emitted
along the backward part is read as absorbed.
\item[(b)] \emph{Energy given to the hits.}  Write \(E=\mu\sinh\chi\), \(p=\mu\cosh\chi\), \(v=\coth\chi\).  If each hit
takes the energy \(\varepsilon\) from the tachyon, \(E\) falls and \(v\) rises.  After \(E_0/\varepsilon\) hits \(E=0\),
\(v=\infty\), \(p=\mu\).  Along the way the cone of r06.2 opens from \(56.68^\circ\) to \(90^\circ\), and the observer
seeing simultaneous hits slows from \(\bsym\) to rest.
\end{enumerate}
In both cases the infinite acceleration is a coordinate effect at a single point.  There the tachyon passes through infinite
speed, rapidity \(i\pi/2\) in the frame at hand.  The proper acceleration stays finite.

\paragraph{Proof sketch.}
(a) \(dx/dt=\coth\sigma\); \(dv/dt=-1/(R\sinh^3\sigma)\); the second derivative with respect to proper length has Minkowski
norm \(1/R\).  (b) \(p^2-E^2=\mu^2\) and \(v=p/E\).  Checked numerically, (b) with six hits from \(v=1/\bsym\).  \(\square\)

\paragraph{Reading.}
A tachyon that pays for the photons of its hits accelerates by itself and ends at infinite speed after finitely many
hits.  At that point the lab frame has become the frame of simultaneous hits, the cone has become a plane wave, and the
tachyon is the pure phase of r06.2(c).  Continuing past this point turns emission into absorption in the lab description.
This is the nearest the hit chain comes to the author's first picture of photons returning to the origin: a chain of
emissions meets a chain of absorptions at the point of infinite speed.

\subsection*{v4.9.0-r06 the Higgs field}
\label{subsec:v490-r06-higgs}

The main body assigns the Higgs sector the role of mediating between the compensated photon fix-sector and the mass
readout of localised branches (Definition~\ref{def:higgs-mediation}, Theorem~\ref{thm:higgs-role-refined}).

\paragraph{Theoremlet v4.9.0-r06.4.}
\begin{enumerate}[leftmargin=2em]
\item[(a)] In the electroweak theory the Higgs field gives the neutral pair \((W^3,B)\) the mass matrix
\(\tfrac{v_H^2}{4}\bigl(\begin{smallmatrix}g^2&-gg'\\-gg'&g'^2\end{smallmatrix}\bigr)\).  It has rank one: the photon is its massless
eigenvector \((\sin\theta_W,\cos\theta_W)\), the \(Z\) the massive one.  This is a port-weighting matrix in the sense of
r05.1, with weight \(0\) on the photon port and \(m_Z^2\) on the other.  For \(g=g'\) (\(\theta_W=45^\circ\)) it is exactly
\(0\cdot C_{\HC}+m_Z^2A\), and the photon is the compensator port.  The measured mixing angle differs from \(45^\circ\), so the
electroweak photon port is a compensator rotated by \(\theta_W-45^\circ\).
\item[(b)] At \(\phi=0\) the Higgs potential has negative mass squared.  The field there is tachyonic in the
field-theoretic sense, and mass arises when it condenses away from \(\phi=0\).
\end{enumerate}

\paragraph{Proof sketch.}
(a) The kernel of the matrix is spanned by \((g',g)\); for \(g=g'\) the matrix is \(\tfrac{v_H^2g^2}{2}A\).  (b) The
potential \(-\mu_H^2|\phi|^2+\lambda|\phi|^4\) has curvature \(-\mu_H^2\) at the origin.  \(\square\)

\paragraph{Reading.}
The Higgs part of the author's picture fits the main body in three respects.
\begin{itemize}[leftmargin=2em]
\item The photon is the direction the Higgs field does not weight, which is Theorem~\ref{thm:higgs-role-refined} in
electroweak form.
\item A field ``trying to create matter'' is a field leaving its tachyonic top.
\item A hit that yields a photon instead of a massive state lands in the unweighted port.
\end{itemize}
Three points need care.
\begin{itemize}[leftmargin=2em]
\item The doublet of r04 is \((E,B)\), the electroweak doublet is \((W^3,B)\).  The correspondence is formal.
\item A field with negative mass squared is an instability that grows; it is not a particle that carries signals faster
than light.  A superluminal chain of hits therefore needs triggers correlated in advance, as for the sweeping light spot,
and not a signal along the line.
\item The condensed Higgs vacuum is the minimum of its potential and supplies no energy.  The energy of the photons comes
either from the tachyon, case r06.3(b), or from the field while it is still rolling off its tachyonic top, that is during
symmetry breaking.
\end{itemize}
Comparing \(\theta_W\) or any other measured quantity with the values of this layer is a prediction and requires a
pre-registration note.

\subsection*{v4.9.0-r06 register entries}
\label{subsec:v490-r06-register}

\begin{center}
\small
\begin{tabular}{@{}>{\raggedright\arraybackslash}p{0.34\textwidth}>{\raggedright\arraybackslash}p{0.58\textwidth}@{}}
\toprule
statement & status after v4.9.0-r06\\
\midrule
r05.3(c), return photon & stands; one photon plus one tachyon leg leaves the gap \(d(1-1/v)\); two tachyon legs can close a loop (r06.1)\\
hit chain at speed \(1/\bsym\) & cone \(56.68^\circ\); all hits simultaneous for the observer at \(\bsym\), where the tachyon is the phase \(i\pi/2\) (r06.2)\\
accelerating tachyon & infinite speed at one point after finite time or finitely many hits; there a pair event, emission turns into absorption (r06.3)\\
Higgs mediation (Def.~\ref{def:higgs-mediation}) & electroweak mass matrix is a port weighting with massless photon port; compensator port for \(\theta_W=45^\circ\); formal correspondence (r06.4)\\
\bottomrule
\end{tabular}
\end{center}

\subsection*{v4.9.0-r06 ledger and guard update}
\label{subsec:v490-r06-ledger-guard}

\[
R_{\rm chain}^{v490\text{-}r06}=(0,0,1)
\]
for (loop condition decided, hit chain and acceleration worked out, physical status of the trigger and of the Higgs
correspondence open).

The layer spends no guard:
\[
R_{\rm glob}^{(v490\text{-}r06,+)}=(0,1,1,1,1),
\qquad
\Gamma_{\rm guard}^{v490\text{-}r06}=(1,1,1,1,1),
\qquad
\beta_4=1/256\ \text{unspent}.
\]
It does not imply exact finite \(C5\) zero closure, exact all-readout terminality,
\(C6/T2\), a global photon-clock/RH theorem, or global all-readout closure.

% ============================================================================
\clearpage
\section*{v4.9.0-r07 Source-Only: The Cone Angle and the Weak Mixing Angle}
\addcontentsline{toc}{section}{v4.9.0-r07 Source-Only: The Cone Angle and the Weak Mixing Angle}
\label{sec:v490-r07-cone-weak-mixing}
% ============================================================================

\paragraph{Purpose.}
This layer records three things.  First, the pre-registered test H14 of the relation between the cone angle of r06.2 and
the weak mixing angle, with its result.  Second, an exploration that followed it, kept as a protocol in which every
variant is counted.  Third, the structural reading that came out of the exploration, and the pre-registered hypothesis
H15 it led to.  Files, all under \texttt{llm/}:
\begin{itemize}[leftmargin=2em]
\item \texttt{Vorab\_Notiz\_H14\_Kegelwinkel\_Weinberg.md};
\item \texttt{Ergebnis\_H14\_Kegelwinkel\_Weinberg.md};
\item \texttt{Erkundungsprotokoll\_Weinberg\_Kegel.md};
\item \texttt{Vorab\_Notiz\_H15\_W\_Masse\_gamma\_L.md};
\item \texttt{Ergebnis\_H15\_W\_Masse\_gamma\_L.md};
\item scripts \texttt{k43\_test\_h14.py} to \texttt{k46\_test\_h15.py}.
\end{itemize}
Data: Particle Data Group, Review of Particle Physics
2024 with the 2025 update (summary tables and the review of the electroweak model), retrieved 2026-10-21.  Earlier layers
are not rewritten.

\noindent\textbf{Source-only status.}
This v4.9.0-r07 layer is source-only.  It does not build a PDF, does not assign a new DOI,
and does not authorize a public release.  The v4.9.0 Version DOI
\href{https://doi.org/10.5281/zenodo.23534954}{10.5281/zenodo.23534954} and the Concept DOI
\href{https://doi.org/10.5281/zenodo.19210728}{10.5281/zenodo.19210728} are unchanged by it.

\subsection*{v4.9.0-r07 test H14}
\label{subsec:v490-r07-h14}

H14 states \(\vartheta_c=2\theta_W\), equivalently \(\sin^2\theta_W=(1-\bsym)/2\).  It was formulated after the closeness
\(56.7^\circ\) against \(2\theta_W\approx57.5^\circ\) had been noticed.  The note says so, fixes the on-shell definition
\(\sin^2\theta_W=1-m_W^2/m_Z^2\) as decisive, and states the expected outcome in advance.

\paragraph{Result.}
The prediction is \(0.22533\).  The on-shell value from the masses is \(0.22321\pm0.00026\), so \(z=+8.2\).  As an exact
relation H14 is \textbf{refuted}.  The prediction lies inside the band spanned by the on-shell and the
\(\overline{\rm MS}\) value (\(0.23122\)), so as a tree-level relation it is \textbf{not decidable}.  A random simple
relation would land in that band with probability of about \(0.02\), so no outcome counts as support.

\subsection*{v4.9.0-r07 exploration, counted}
\label{subsec:v490-r07-exploration}

The exploration asked whether a physical effect can close the gap of \(+0.00212\) in \(\sin^2\theta_W\).  The gap is
equivalent to removing a contribution \(\Delta\rho_X=+0.00275\) present in nature, or to an \(m_W\) smaller by
\(110\) MeV.  The variants tried, each counted, were:
\begin{itemize}[leftmargin=2em]
\item the tribonacci value in place of \(\bsym\), which moves the prediction the wrong way;
\item removal of the top-quark contribution, which overshoots to \(0.2304\);
\item the tree-level angle from \((\alpha,G_F,m_Z)\), giving \(0.212\) or \(0.234\) depending on \(\alpha\);
\item a fraction of the top contribution, with no meaning;
\item neutrinos, eleven orders of magnitude too small;
\item every combination of omitted contributions, of which only two are large (top \(0.00934\), remainder \(0.00108\))
and none hits;
\item mixing with a dark state.
\end{itemize}
The last variant gives a general bound.  Any real contribution of the required size would shift \(m_W\) by \(155\) MeV
against the Standard-Model prediction.  The measured value differs from that prediction by \(+10\pm14\) MeV.  Real effects
of particles, known or dark, can therefore account for at most about a quarter of the gap.

\subsection*{v4.9.0-r07 the structural reading}
\label{subsec:v490-r07-structure}

\paragraph{Theoremlet v4.9.0-r07.1.}
With \(\Pi(x)=2\arctan e^{-x}\), Lobachevsky's angle of parallelism, the following are equivalent:
\[
\cos2\theta=\tanh\eta,\qquad \tan\theta=e^{-\eta},\qquad \sin^2\theta=\frac1{1+e^{2\eta}},\qquad 2\theta=\Pi(\eta).
\]
The cone angle of r06.2 is \(\vartheta_c=\Pi(\eta)\).  For the port weighting \(e^{t\sigma_x}=e^tC_{\HC}+e^{-t}A\) of
r05.1(d) at \(t=\eta\), the ratio of the weights is \(e^{2\eta}=\cot^2\theta\).

\paragraph{Proof sketch.}
\(\cos2\theta=(1-\tan^2\theta)/(1+\tan^2\theta)\) with \(\tan^2\theta=e^{-2\eta}\) gives \(\tanh\eta\);
\(\cos\Pi(x)=\tanh x\).  Checked numerically.  \(\square\)

\paragraph{Consequence.}
H14 is a statement about a ratio of couplings, \(g/g'=e^{\eta}\), the Doppler factor of way (C) in r01.  It is not a
statement about a ratio of masses.  Couplings run.  In the \(\overline{\rm MS}\) scheme \(\sin^2\theta_W\) has its minimum,
about \(0.231\), at \(\mu=M_W\), and rises above and below.  No scale reaches \(0.22533\).  The structural reading thus
points to the scheme in which the gap is larger, and it does not support the closeness to the on-shell value.

\subsection*{v4.9.0-r07 an observation, and hypothesis H15}
\label{subsec:v490-r07-h15}

\paragraph{Observation (found during the exploration).}
The on-shell coupling ratio is \(\cot\theta_W=m_W/\sqrt{m_Z^2-m_W^2}=1.86550\pm0.00140\).  The three constants of the
source in this range sit at the following distances: \(\gL=\cosh2\eta=1.86444\) at \(z=-0.8\), \(e^{\eta}=1.85419\) at
\(z=-8.1\), \(1/\bsym=1.82034\) at \(z=-32\).  By Theoremlet r01.3, \(\gL\) (way A) and \(e^\eta\) (way C) coincide at the
tribonacci point and differ by \(0.55\%\) otherwise.  The data sit at way (A).  The cone relation used way (C).

\paragraph{Why this is not evidence.}
The relation was found after the fact.  The three constants \(1.820\), \(1.854\), \(1.864\) lie in the range that
\(\cot\theta_W\) covers across schemes, \(1.823\) to \(1.866\), so some pairing is nearly expected.  The chance of a
\(1\sigma\) hit of one of the three at the on-shell value is about \(0.17\).  Reading~B contains no
\(SU(2)\times U(1)\), so the identification of the weight ratio with \((g/g')^2\) is formal.

\paragraph{Hypothesis H15 (pre-registered).}
\(\cot\theta_W=\gL\) in the on-shell definition, that is \(m_W=m_Z\gL/\sqrt{1+\gL^2}\).  Here \(\gBI\) is the root of the
state-counting equation \(\sum_{j}2\exp(-2\pi\gamma\sqrt{j(j+1)})=1\) that r03 cites, and \(\gstring=\ln2/(\pi\sqrt3)\).
Computed immediately after the note was committed and before any comparison:
\[
\gBI=0.237532957966,\qquad \gL=1.8646997646,\qquad m_W^{\rm pred}=80.3615\pm0.0018\ \text{GeV}.
\]
The comparison is with the first world average or combination of \(m_W\) published after 2026-10-21.  Known values do not
count.  The prediction lies within a few MeV of the Standard-Model prediction (\(80.359\pm0.006\) GeV).  A comparison can
therefore refute H15, but it can support H15 over the Standard Model only once the measurement excludes the
Standard-Model value at \(3\sigma\), which needs a precision below about \(1\) MeV.  The result is pending.

\subsection*{v4.9.0-r07 register entries}
\label{subsec:v490-r07-register}

\begin{center}
\small
\begin{tabular}{@{}>{\raggedright\arraybackslash}p{0.34\textwidth}>{\raggedright\arraybackslash}p{0.58\textwidth}@{}}
\toprule
statement & status after v4.9.0-r07\\
\midrule
H14, \(\vartheta_c=2\theta_W\) & refuted as exact relation (\(z=+8.2\)); not decidable at tree level; no support\\
cone angle & \(\vartheta_c=\Pi(\eta)\), the angle of parallelism (r07.1); not measured\\
gap of H14 by real particle effects & bounded by \(m_W\) to at most about a quarter (exploration, 8 counted variants)\\
H15, \(\cot\theta_W=\gL\) & pre-registered; prediction \(m_W=80.3615\pm0.0018\) GeV; comparison pending\\
r01 ways (A)/(C) & the on-shell coupling ratio lies at way (A) (observation, not a test)\\
\bottomrule
\end{tabular}
\end{center}

\subsection*{v4.9.0-r07 ledger and guard update}
\label{subsec:v490-r07-ledger-guard}

\[
R_{\rm ew}^{v490\text{-}r07}=(0,0,1)
\]
for (H14 decided, structural reading of the cone established, H15 pending).

The layer spends no guard:
\[
R_{\rm glob}^{(v490\text{-}r07,+)}=(0,1,1,1,1),
\qquad
\Gamma_{\rm guard}^{v490\text{-}r07}=(1,1,1,1,1),
\qquad
\beta_4=1/256\ \text{unspent}.
\]
It does not imply exact finite \(C5\) zero closure, exact all-readout terminality,
\(C6/T2\), a global photon-clock/RH theorem, or global all-readout closure.

% ============================================================================
\clearpage
\section*{v4.9.0-r08 Source-Only: Frame Drift, Atomic Clocks and the Hubble Tension}
\addcontentsline{toc}{section}{v4.9.0-r08 Source-Only: Frame Drift, Atomic Clocks and the Hubble Tension}
\label{sec:v490-r08-drift-clocks-hubble}
% ============================================================================

\paragraph{Purpose.}
Proposition~\ref{prop:H0-tension} explains the Hubble tension by a drift of the frame offset.  r17 and r18 tested that
drift against the expansion history.  This layer tests it against the constancy of the fine-structure constant, using
the source's own rule for which quantities stay fixed.  It also checks the other explanations of the tension available
in the source, settles the sign of \(k\) for \(\alpha\) left open in r03, names a frame-invariant clock, and states what
the drift means for H15 of r07.  Files: \texttt{llm/Erkundungsprotokoll\_Weinberg\_Kegel.md} (entries E10 to E13) and
scripts \texttt{llm/k47\_drift\_weinberg.py}, \texttt{llm/k48\_feste\_groessen.py}.  No new prediction is made.  Earlier
layers are not rewritten.

\noindent\textbf{Source-only status.}
This v4.9.0-r08 layer is source-only.  It does not build a PDF, does not assign a new DOI,
and does not authorize a public release.  The v4.9.0 Version DOI
\href{https://doi.org/10.5281/zenodo.23534954}{10.5281/zenodo.23534954} and the Concept DOI
\href{https://doi.org/10.5281/zenodo.19210728}{10.5281/zenodo.19210728} are unchanged by it.

\subsection*{v4.9.0-r08 what stays fixed under drift}
\label{subsec:v490-r08-fixed}

The source fixes this in four places.
\begin{itemize}[leftmargin=2em]
\item The definition in \S\ref{sec:phase-constants}: \(X^{(\mathrm{photon})}=\gL^{-k(X)}X^{(\mathrm{matter})}\).
\item Theorem~\ref{thm:k-values}: \(k(c)=0\), \(k(\gBI)=0\) (``frame-invariant by construction'').
\item The corollary ``\(\alpha\)-drift as a frame effect'' in the same section: ``the hardware value
\(\alpha^{(\mathrm{photon})}\) is strictly constant'', with \(\alpha^{(\mathrm{matter})}(t)=\gL^{(\mathrm{local})}(t)\,
\alpha^{(\mathrm{photon})}\).
\item Item (1) of Proposition~\ref{prop:saturation}, which writes the neutrino mass in the same way.
\end{itemize}
Photon-frame values and quantities with \(k=0\) stay fixed.  Matter-frame values drift with a local factor
\(\gL^{(\mathrm{local})}(t)^{k}\).  The structural factor \(\gL=\gBI/\gstring\) is fixed, since both inputs are fixed
numbers.  The source thus separates a fixed structural factor from a local one that may drift.  This is the separation
r18 offered.

\subsection*{v4.9.0-r08 the sign of \texorpdfstring{$k$}{k} for \texorpdfstring{$\alpha$}{alpha}}
\label{subsec:v490-r08-alpha-sign}

r03 (F2) found that Theorem~\ref{thm:k-values}, whose row reads ``\(\alpha^{-1}\ldots k=1\)'', and the corollary above are
inverse to each other, and left the choice open.  The main body uses \(k\) for \(\alpha\) in three places:
\begin{itemize}[leftmargin=2em]
\item prediction P5 among the predictions around Corollary~\ref{cor:falsifiable}, with
\(\gL^{k(\alpha)-k(G_N)}=\gL^{-1}\);
\item Proposition~\ref{prop:muon-g2}, ``\(k(\alpha_s)=1\) (like the electromagnetic coupling \(\alpha\))'' with
\(\alpha_s^{(\mathrm{photon})}=\gL^{-1}\alpha_s^{(\mathrm{matter})}\);
\item the corollary, \(\alpha^{(\mathrm{matter})}=\gL\,\alpha^{(\mathrm{photon})}\).
\end{itemize}
All three mean \(k(\alpha)=+1\) for \(\alpha\) itself.  No photon-frame value of \(\alpha^{-1}\) is used anywhere in the main
body.

\paragraph{Resolution (reading decision).}
The table row shows the CODATA value of \(\alpha^{-1}\), and its entry \(k=1\) belongs to \(\alpha\).  With this reading all
three uses agree and nothing else changes.  The resolution changes no number of the main body.

\subsection*{v4.9.0-r08 the drift against atomic clocks}
\label{subsec:v490-r08-clocks}

\paragraph{Theoremlet v4.9.0-r08.1.}
Suppose the drift of Proposition~\ref{prop:H0-tension}, \(1.6\times10^{-12}\) per year in \(\bsym\), acts through the local
frame factor, as r18 reads it.  Then the corollary gives \(|{\rm d}\ln\alpha/{\rm d}t|\) for each choice of the factor:

\begin{center}
\small
\begin{tabular}{@{}lrrr@{}}
\toprule
factor & exponent needed for \(1.083\) & \(|{\rm d}\ln\alpha/{\rm d}t|\) per year & \(|\Delta\alpha/\alpha|\) since \(z=1100\)\\
\midrule
\(\gL=\cosh2\eta\) & \(1.52\) & \(3.9\times10^{-12}\) & \(5.4\%\)\\
neighbour, \(\cosh\eta\) & \(4.72\) & \(1.3\times10^{-12}\) & \(1.7\%\)\\
neighbour, \(e^{\eta}\) & \(2.54\) & \(2.3\times10^{-12}\) & \(3.2\%\)\\
\bottomrule
\end{tabular}
\end{center}

The comparison of two \({}^{171}\mathrm{Yb}^+\) clocks gives \((1/\alpha)\,{\rm d}\alpha/{\rm d}t=(1.0\pm1.1)\times10^{-18}\)
per year (Lange et al., Phys.\ Rev.\ Lett.\ 126, 011102, 2021).  Every choice lies about six orders of magnitude above
that bound.

Independently of the reading: if a drifting factor with exponent \(k(H_0)\) produces the tension \(\ln1.083\), then
\(\alpha\) changes by \(\ln1.083\cdot k(\alpha)/k(H_0)\).  For a drift linear in time, the clocks require
\(k(\alpha)/k(H_0)<5.5\times10^{-7}\) at \(2\sigma\).

\paragraph{Proof sketch.}
\({\rm d}\ln\gL/{\rm d}t=4\bsym(1-\bsym^2)^{-2}\dot\bsym/\gL\), and correspondingly for the neighbour factors.  The
exponents follow from the ratios of the factors at \(z=0\) and \(z=1100\).  Checked numerically.  \(\square\)

\paragraph{Why measuring through photons does not help.}
A clock comparison measures the ratio of two transition frequencies of atoms.  A common frame factor on both frequencies
cancels.  The dependence on \(\alpha\) comes from inside the atom, which is matter.  If the atom's internal physics were
photon-frame physics, there would be no drifting \(\alpha^{(\mathrm{matter})}\) at all.  The corollary's annual-drift
prediction would then lapse too, and Proposition~\ref{prop:H0-tension} would lose its mechanism.

\paragraph{Consequence.}
In its frame-factor form, Proposition~\ref{prop:H0-tension} cannot explain the tension and keep \(\alpha\) constant at the
same time.  Two ways out would remain, and the source contains neither.  One is \(k(\alpha)\approx0\), which contradicts the
resolution above and the corollary.  The other is a separate factor for the expansion that does not touch the
constants.  r03 already bounded the annual modulation of the local factor by \(10^{-17}\).  The present layer applies the
same clocks to the secular drift.

\subsection*{v4.9.0-r08 other explanations in the source}
\label{subsec:v490-r08-other}

\begin{enumerate}[leftmargin=2em]
\item \emph{Static frame typing.}  The value inferred from the microwave background would be read in another frame than
the local value.  A constant factor disturbs no clock.  By Corollary~\ref{cor:gue-frame} a clock frequency transforms with
\(\gL\), so \(k(H_0)=1\), which gives a factor \(1.864\) (neighbour factors \(1.197\), \(1.854\)) against the observed
\(1.083\pm0.017\).  Only fractional exponents near \(0.13\) or \(0.44\) would fit, and Theorem~\ref{thm:k-values} gives no
basis for them.  The observed ratio is too loose to single out an exponent.
\item \emph{A change confined to the early universe.}  A larger electron mass at recombination, by about one to two
percent, is discussed in the literature as a relief of the tension.  With the drift of
Proposition~\ref{prop:H0-tension}, \(\bsym\) was smaller earlier, so \(m_e^{(\mathrm{matter})}\) with \(k=1\) was smaller.
That is the wrong direction, and the same holds for \(\alpha\) with \(k(\alpha)=+1\).  The opposite sign of the drift, the one
r17 found Theorem~\ref{thm:dark-energy} to require, would point the right way.  The change would then have to stop before
\(z\approx3\), because of quasar spectra and clocks, and at \(5\%\) it would be too large.
\end{enumerate}
No explanation of the tension within the source survives the clocks, the quasar spectra and the size of the effect
together.  Proposition~\ref{prop:H0-tension} stays a hypothesis without a working mechanism.

\subsection*{v4.9.0-r08 a frame-invariant clock}
\label{subsec:v490-r08-compton}

\paragraph{Theoremlet v4.9.0-r08.2.}
By Theorem~\ref{thm:k-values} the Compton frequency \(m_ec^2/\hbar\) has \(k=k(m_ec^2)-k(\hbar)=1-1=0\).  A clock ticking
at a Compton frequency is frame-invariant in the source's sense.  An atomic transition frequency scales as
\(\alpha^2m_ec^2/\hbar\) and has \(k=2k(\alpha)=2\).  The ratio of the two clocks reads \(\gL^{(\mathrm{local})}(t)^{\pm2}\)
directly, a local test of the frame factor that needs no cosmology.

\paragraph{Proof sketch.}
Additivity of \(k\) (Lemma v4.8.0-r03.1).  \(\square\)

\paragraph{Remark.}
Clocks tied to the Compton frequency of a mass have been built by atom interferometry (Lan et al., Science 339, 554, 2013).
A comparison of such a clock with an atomic clock would read the drift of r08.1 at twice the rate.  Existing clock
comparisons already exclude that drift.  The clock concerns local readout only.  It touches neither the local selector
guard nor the forbidden global photon-clock statement.

\subsection*{v4.9.0-r08 consequence for H15}
\label{subsec:v490-r08-h15}

H15 (r07) is stated with the structural factor \(\gL=\gBI/\gstring\), which is fixed.  It predicts a weak mixing angle that
is constant in time.  In addition, \(m_W/m_Z\) is a ratio of masses, and if both carry the same \(k\) the frame factor
cancels.  A drift of \(\theta_W\), with \(\sin^2\theta_W=0.242\) at \(z=1100\), would follow only if H15 were read with the
local factor, which the source's rules do not suggest.  The pre-registered comparison of H15 is unaffected.

\subsection*{v4.9.0-r08 register entries}
\label{subsec:v490-r08-register}

\begin{center}
\small
\begin{tabular}{@{}>{\raggedright\arraybackslash}p{0.34\textwidth}>{\raggedright\arraybackslash}p{0.58\textwidth}@{}}
\toprule
statement & status after v4.9.0-r08\\
\midrule
Prop.~\ref{prop:H0-tension} & frame-factor drift incompatible with atomic clocks by about six orders of magnitude in every reading (r08.1); static typing and early-only change do not work; hypothesis without working mechanism\\
Thm.~\ref{thm:k-values}, row \(\alpha^{-1}\) & resolved: \(k=1\) belongs to \(\alpha\); consistent with P5, Prop.~\ref{prop:muon-g2} and the corollary (closes r03 F2)\\
fixed quantities under drift & photon-frame values and \(k=0\); structural \(\gL=\gBI/\gstring\) fixed; local factor may drift\\
Compton clock & \(k=0\), frame-invariant; ratio to an atomic clock reads \(\gL^{(\mathrm{local})}{}^{\pm2}\) (r08.2)\\
H15 & predicts a constant \(\theta_W\); comparison unaffected\\
\bottomrule
\end{tabular}
\end{center}

\subsection*{v4.9.0-r08 ledger and guard update}
\label{subsec:v490-r08-ledger-guard}

\[
R_{\rm drift}^{v490\text{-}r08}=(0,0,1)
\]
for (drift against clocks decided, sign of \(k(\alpha)\) resolved, a mechanism for the Hubble tension open).

The layer spends no guard:
\[
R_{\rm glob}^{(v490\text{-}r08,+)}=(0,1,1,1,1),
\qquad
\Gamma_{\rm guard}^{v490\text{-}r08}=(1,1,1,1,1),
\qquad
\beta_4=1/256\ \text{unspent}.
\]
It does not imply exact finite \(C5\) zero closure, exact all-readout terminality,
\(C6/T2\), a global photon-clock/RH theorem, or global all-readout closure.

% ============================================================================
\clearpage
\section*{v4.9.0-r09 Source-Only: What the Local Frame Factor Can Be}
\addcontentsline{toc}{section}{v4.9.0-r09 Source-Only: What the Local Frame Factor Can Be}
\label{sec:v490-r09-local-frame-factor}
% ============================================================================

\paragraph{Purpose.}
r08 separated a fixed structural factor \(\gL=\gBI/\gstring\) from a local factor \(\gL^{(\mathrm{local})}\) that may vary.
This layer collects what the source says the local factor depends on.  It records two contradictions found on the way,
tests three readings of the local factor against clock data, and corrects the arithmetic of prediction P5.  Files:
\texttt{llm/Erkundungsprotokoll\_Weinberg\_Kegel.md} (entries E14, E15) and the script
\texttt{llm/k49\_lokaler\_faktor.py}.  No new prediction is made.  Earlier layers are not rewritten.

\noindent\textbf{Source-only status.}
This v4.9.0-r09 layer is source-only.  It does not build a PDF, does not assign a new DOI,
and does not authorize a public release.  The v4.9.0 Version DOI
\href{https://doi.org/10.5281/zenodo.23534954}{10.5281/zenodo.23534954} and the Concept DOI
\href{https://doi.org/10.5281/zenodo.19210728}{10.5281/zenodo.19210728} are unchanged by it.

\subsection*{v4.9.0-r09 what the source says}
\label{subsec:v490-r09-source}

The source names four influences on the local factor.  For none of them does it give a derived relation, and it gives
no field law for \(\gL^{(\mathrm{local})}(x,t)\).
\begin{center}
\small
\begin{tabular}{@{}>{\raggedright\arraybackslash}p{0.17\textwidth}>{\raggedright\arraybackslash}p{0.33\textwidth}>{\raggedright\arraybackslash}p{0.42\textwidth}@{}}
\toprule
influence & place & content\\
\midrule
observer & Theorem~\ref{thm:frame-offset} & every matter observer carries \(\zminus\neq0\); measured constants are frame-invariant values times a Lorentz factor of the observer\\
cosmic time & Proposition~\ref{prop:H0-tension} & \(\bsym\) grows ``as matter structures in the universe grew''\\
Earth's orbit & corollary ``\(\alpha\)-drift as a frame effect'' & the factor varies ``as Earth moves through the Sphaera-Cascade geometry over its orbit'', with a period of one year, absent for a clock that does not co-orbit\\
potential & prediction P5 near Corollary~\ref{cor:falsifiable} & amplitude of the \(\alpha\) modulation \(\sim10^{-10}\,\delta\Phi/c^2\)\\
\bottomrule
\end{tabular}
\end{center}

\paragraph{Two contradictions.}
\begin{enumerate}[leftmargin=2em]
\item Theorem~\ref{thm:frame-offset} calls \(\gBI\) ``the Lorentz-transformed \(\gstring\), observed from a matter frame''.
In that theorem \(\gBI\) depends on the observer.  Theorem~\ref{thm:k-values} sets \(k(\gBI)=0\), ``frame-invariant by
construction''.
\item Theorem~\ref{thm:frame-offset} projects all constants, including \(c\), with one Lorentz factor, which means
\(k=1\) for all.  Theorem~\ref{thm:k-values} gives \(k(c)=0\) and \(k(G_N)=2\).
\end{enumerate}
Both are resolved if Theorem~\ref{thm:frame-offset} is read as a statement about the structural factor alone.  That is the
separation of r18 and r08.  What the local factor consists of then remains open.

\subsection*{v4.9.0-r09 three readings tested}
\label{subsec:v490-r09-readings}

With the rule \(\delta\ln X=k(X)\,\delta\ln\gL^{(\mathrm{local})}\), three readings were tested.

\paragraph{Theoremlet v4.9.0-r09.1.}
\begin{enumerate}[leftmargin=2em]
\item[(V)] \emph{Velocity against a preferred frame.}  If the observer's rapidity adds to \(\eta\), Earth's orbital motion
projected on the dipole of the microwave background gives \(\delta\eta\approx9.8\times10^{-5}\).  That yields
\(\delta\ln\gL=2\tanh2\eta\,\delta\eta\approx1.6\times10^{-4}\), and with \(k(\alpha)=1\) an annual \(\alpha\) modulation of the
same size.  The clock fit quoted in r03 (F2) allows about \(2\times10^{-17}\).  The reading is \textbf{excluded}, by a factor
of about \(10^{13}\).  The static part, \(370\) km/s against the background, would make \(\alpha\) depend on direction at the
level of \(2\times10^{-3}\).
\item[(P)] \emph{Gravitational potential.}  \(\delta\ln\gL^{(\mathrm{local})}=\kappa\,\delta\Phi/c^2\).  The source's own
coefficient \(\kappa\sim10^{-10}\) and the eccentricity of Earth's orbit (\(\delta\Phi/c^2\approx1.6\times10^{-10}\)) give an
annual \(\delta\alpha/\alpha\approx1.6\times10^{-20}\), inside the bounds.  The reading \textbf{survives}.  It cannot drive the
cosmological drift of Proposition~\ref{prop:H0-tension}: potentials of large-scale structure, \(\sim10^{-5}\), give
\(\delta\ln\gL\sim10^{-15}\), against about \(5\times10^{-2}\) needed.
\item[(D)] \emph{Matter density.}  The source gives no functional form.  Density differs between laboratory, space and
astrophysical environments by many orders of magnitude, so any appreciable coupling would have to be screened.  The
reading is \textbf{not testable} as stated.
\end{enumerate}

\paragraph{Proof sketch.}
(V) \({\rm d}\ln\cosh2\eta/{\rm d}\eta=2\tanh2\eta\), with \(29.8\cos11^\circ\) km/s against \(c\).  (P) \(GM_\odot/(c^2a)=9.87\times
10^{-9}\) times the eccentricity \(0.0167\).  Checked numerically.  \(\square\)

\paragraph{Consequence.}
Of the three readings only the potential reading is compatible with the clocks.  It is also the one the source itself
states in P5.  It describes a small violation of local position invariance and cannot explain the Hubble tension.
Within it, every constant couples to the potential with \(\beta_X=k(X)\,\kappa\):
\begin{itemize}[leftmargin=2em]
\item the Compton clock of r08.2, with \(k=0\), does not couple;
\item \(\alpha\) couples with \(\kappa\);
\item an atomic clock and \(G_N\), with \(k=2\), couple with \(2\kappa\).
\end{itemize}
The ratio of the couplings, not their size, is the content of the source.  The size \(\kappa\) is not derived.

\subsection*{v4.9.0-r09 the arithmetic of P5}
\label{subsec:v490-r09-p5}

P5 states
\(\frac{\delta\alpha/\alpha}{\delta G_N/G_N}\big|_{\rm annual}=\gL^{k(\alpha)-k(G_N)}\cdot\frac{k(\alpha)}{k(G_N)}=\gL^{-1}\approx
0.537\).

\paragraph{Theoremlet v4.9.0-r09.2.}
From \(X^{(\mathrm{matter})}=\gL^{k}X^{(\mathrm{photon})}\) with fixed photon-frame values follows
\(\delta X/X=k(X)\,\delta\gL/\gL\).  Hence
\[
\frac{\delta\alpha/\alpha}{\delta G_N/G_N}=\frac{k(\alpha)}{k(G_N)}=\frac12 .
\]
The static factor \(\gL^{k(\alpha)-k(G_N)}\) does not belong in a ratio of relative changes.  The middle expression of P5
equals \(\gL^{-1}/2=0.268\), not \(\gL^{-1}\).  The corrected value \(0.5\) is the same in every reading of the local factor.

\paragraph{Proof sketch.}
Logarithmic differentiation.  \(\square\)

\subsection*{v4.9.0-r09 register entries}
\label{subsec:v490-r09-register}

\begin{center}
\small
\begin{tabular}{@{}>{\raggedright\arraybackslash}p{0.34\textwidth}>{\raggedright\arraybackslash}p{0.58\textwidth}@{}}
\toprule
statement & status after v4.9.0-r09\\
\midrule
local factor \(\gL^{(\mathrm{local})}\) & four influences named (observer, time, orbit, potential), no field law; only the potential reading survives the clocks (r09.1)\\
Thm.~\ref{thm:frame-offset} vs Thm.~\ref{thm:k-values} & two contradictions (\(\gBI\) observer-dependent vs \(k=0\); one factor for all vs \(k(c)=0\), \(k(G_N)=2\)); resolved by reading Thm.~\ref{thm:frame-offset} as structural\\
P5, annual ratio & corrected to \(k(\alpha)/k(G_N)=1/2\); the stated \(\gL^{-1}=0.537\) is wrong twice (r09.2)\\
potential coupling & \(\beta_X=k(X)\kappa\); Compton clock uncoupled; \(\kappa\sim10^{-10}\) not derived, below present bounds\\
Prop.~\ref{prop:H0-tension} & the surviving reading gives \(10^{-15}\) instead of \(5\times10^{-2}\); unchanged: hypothesis without working mechanism\\
\bottomrule
\end{tabular}
\end{center}

\subsection*{v4.9.0-r09 ledger and guard update}
\label{subsec:v490-r09-ledger-guard}

\[
R_{\rm loc}^{v490\text{-}r09}=(0,0,1)
\]
for (influences and contradictions recorded, readings of the local factor decided against clocks, size of \(\kappa\) open).

The layer spends no guard:
\[
R_{\rm glob}^{(v490\text{-}r09,+)}=(0,1,1,1,1),
\qquad
\Gamma_{\rm guard}^{v490\text{-}r09}=(1,1,1,1,1),
\qquad
\beta_4=1/256\ \text{unspent}.
\]
It does not imply exact finite \(C5\) zero closure, exact all-readout terminality,
\(C6/T2\), a global photon-clock/RH theorem, or global all-readout closure.

% ============================================================================
\clearpage
\section*{v4.9.0-r10 Source-Only: A Reading of the C5 Factor 3/4 and a Class of Shape-(iii) Witnesses}
\addcontentsline{toc}{section}{v4.9.0-r10 Source-Only: A Reading of the C5 Factor 3/4 and a Class of Shape-(iii) Witnesses}
\label{sec:v490-r10-c5-factor-shape-three}
% ============================================================================

\paragraph{Purpose.}
r94 found that the C5 contraction factor \(3/4\) is declared, not derived, and that a single projection explains at most
one step.  r88 left shape (iii), a witness outside the collar iteration, as the only open route to a finite C5 discharge.
This layer offers a reading of the factor built from objects the source already has, and uses the same reading to
exclude a class of shape-(iii) witnesses.  The reading is an interpretation, admitted as such by the author
(2026-10-22), and is not a derivation.  Files: \texttt{llm/Erkundungsprotokoll\_Kern.md} and the scripts
\texttt{llm/k50\_c5\_faktor.py}, \texttt{llm/k51\_c5\_form\_iii.py}.  Earlier layers are not rewritten.
This layer does not assert or approach an exact finite C5 zero closure.

\noindent\textbf{Source-only status.}
This v4.9.0-r10 layer is source-only.  It does not build a PDF, does not assign a new DOI,
and does not authorize a public release.  The v4.9.0 Version DOI
\href{https://doi.org/10.5281/zenodo.23534954}{10.5281/zenodo.23534954} and the Concept DOI
\href{https://doi.org/10.5281/zenodo.19210728}{10.5281/zenodo.19210728} are unchanged by it.

\subsection*{v4.9.0-r10 rotate, then compensate}
\label{subsec:v490-r10-rotate-compensate}

On the doublet of r04, with \(C_{\HC}=\tfrac12(1+\sigma_x)\), \(A=\tfrac12(1-\sigma_x)\) and the duality rotation
\(R_\alpha\), write \(u\) for the unit vector spanning \(\operatorname{ran}C_{\HC}\).

\paragraph{Theoremlet v4.9.0-r10.1.}
\(C_{\HC}R_\alpha C_{\HC}=\cos\alpha\,C_{\HC}\).  One step ``rotate by \(\alpha\), then compensate'' keeps the fraction
\(\cos^2\alpha\) of the energy in the compensator port and sends \(\sin^2\alpha\) to the partner port.  Because the rotation
lies between the projections, the step is not idempotent: every repetition removes the same fraction again.  For
\(\alpha=\pi/6\):
\[
\cos^2\tfrac\pi6=\tfrac34,\qquad
\bigl(\tfrac34\bigr)^n:\ \tfrac34,\ \tfrac9{16},\ \tfrac{27}{64},\ \tfrac{81}{256},\ \ldots,
\]
which are the ledger values of r04, r05, r07 and r08.  The first stage with \((3/4)^n\le1/256\) is \(n=20\).

\paragraph{Proof sketch.}
\(\langle u,R_\alpha u\rangle=\cos\alpha\) in two dimensions.  The stage was checked in exact rational arithmetic.
\(\square\)

\paragraph{Reading.}
The step meets the condition r94 set for any derivation: a refreshed one-quarter identity at every stage.  Two sources for
\(\pi/6\) are available in the source.
\begin{itemize}[leftmargin=2em]
\item \emph{Triolic.}  The quarter turn from matter to tachyon (\(\pi/2\), Theoremlet r07.2) is divided into three equal
steps.  Then \(27/64\), the r07 probe entry, is one complete passage from matter to tachyon.  This reading needs no value of
\(\bsym\).
\item \emph{Hopf.}  \(\tan\alpha=1/\sqrt3\), the Hopf value of \(\bsym\) in Theorem~\ref{thm:dark-energy}, so
\(q=1/(1+\bsym^2)=3/4\).  With the measured \(\bsym=0.5493\) the same formula gives \(q=0.7682\) and stage \(22\), outside the
interval of Theoremlet r94.2.  This reading holds only at the Hopf value.
\end{itemize}

\paragraph{A Zeno pattern.}
A single quarter turn followed by compensation empties the compensator port, \(C_{\HC}R_{\pi/2}u=0\).  Three compensated
steps of \(\pi/6\) leave \(27/64\), six leave \(0.66\), three hundred leave \(0.99\).  The more often the content is
compensated, the less of it leaves.  This is the pattern of the quantum Zeno effect.

\paragraph{Observations, not interpretations.}
\begin{enumerate}[leftmargin=2em]
\item The relay couplings of the r14 chart family take the values \(0,\tfrac14,\tfrac12,\tfrac34,1\).  These are
\(\sin^2\) of \(0,\pi/6,\pi/4,\pi/3,\pi/2\).  A uniform grid of quarters has this property anyway, so it carries little
weight.
\item Chart \(V3\), whose relay coupling \(\kappa=1/4\) produces the five in \(\Delta_\gamma<1/5\) (r92), and the clock seed,
\((N,\kappa,G)=(5,\tfrac34,1)\), have couplings that are the two port fractions of one \(\pi/6\) step.  In \(V3\),
\(\kappa(N-1)=1\) doubles the overlap flux, \(S=2N=10\), and gives \(G^{\rm req}/G^{\rm avail}=4/5\).
\item The tolerated collar row of r05 has ratio \(5/6=1/(1+\tfrac15)\), the form of the Hopf reading with \(1/5\) in place of
\(1/3\).
\end{enumerate}
None of these is used below.

\subsection*{v4.9.0-r10 a class of shape-(iii) witnesses}
\label{subsec:v490-r10-shape-three}

r88 requires of a shape-(iii) witness that it exhibit the vanishing of the obstruction itself, not of the score
\((\alpha)\), and that it not import the all-readout slot \((\delta)\).  Under the reading above, the C5 content is the field
content of the doublet, and the score is its energy in the compensator port.  The transports available in Reading~B are
spatial rotations, duality rotations and boosts.

\paragraph{Theoremlet v4.9.0-r10.2 (conditional on the reading).}
\begin{enumerate}[leftmargin=2em]
\item[(a)] \emph{Isometric witnesses} (spatial and duality rotations) conserve
\(\|FC_{\HC}\|^2+\|FA\|^2=\|F\|^2\).  A witness of this kind that empties the compensator port puts the whole content into the
partner port.  It satisfies \((\alpha)\) only by importing the partner readout, which is \((\delta)\).
\item[(b)] \emph{Boosts} change the total \(E^2+B^2\), but \(E^2+B^2\ge|s|\) for every boost, with equality in the frame
where \(E\parallel B\).  For a non-null field the content therefore has a positive floor that no boost removes.
\item[(c)] For a null field the total falls like \(e^{-2\eta}\) under a boost along the propagation, and reaches zero only in
the limit \(\eta\to\infty\), never at a finite stage.
\end{enumerate}
Hence no witness built from the transports of Reading~B meets \((\alpha)\) and \((\delta)\) at a finite stage.  The obstruction
behaves like the invariant \(s\), which no group action of Reading~B changes (Theoremlet r19.1).

\paragraph{Proof sketch.}
(a) \(C_{\HC}+A=1\), \(C_{\HC}A=0\), both symmetric.  (b) \(u^2-|E\times B|^2=\tfrac14|s|^2\) (r03.3) gives \(2u\ge|s|\); checked
over \(16200\) boosts.  (c) Doppler factor of a plane wave.  \(\square\)

\paragraph{Consequence.}
Under this reading, a shape-(iii) witness would have to lie outside the transport group of Reading~B.  It would have to
change the invariant \(s\), that is, carry dynamics.  Reading~B has no dynamics (r10 of v4.8.0).  The route therefore stays
open in the sense of r88, and the class of witnesses examined here is closed.  The layer does not assert, and does not
approach, an exact finite C5 zero closure.  It records why the transports at hand cannot supply one.

\subsection*{v4.9.0-r10 register entries}
\label{subsec:v490-r10-register}

\begin{center}
\small
\begin{tabular}{@{}>{\raggedright\arraybackslash}p{0.34\textwidth}>{\raggedright\arraybackslash}p{0.58\textwidth}@{}}
\toprule
statement & status after v4.9.0-r10\\
\midrule
C5 factor \(3/4\) (r94) & reading offered: rotate by \(\pi/6\), then compensate; reproduces \(3/4,9/16,27/64,81/256\) and \(n=20\); interpretation, not derivation (r10.1)\\
source of \(\pi/6\) & triolic division of the matter--tachyon quarter turn, or Hopf value \(1/\sqrt3\) (only at that value)\\
shape (iii) (r88) & under the reading: isometric witnesses relocate, boosts meet the floor \(|s|\), null fields reach zero only as a limit; witnesses would need dynamics (r10.2)\\
\(\Delta_\gamma<1/5\) (r92) & observation: \(V3\) coupling \(1/4\) and clock-seed coupling \(3/4\) are the two port fractions of one \(\pi/6\) step; not used\\
\bottomrule
\end{tabular}
\end{center}

\subsection*{v4.9.0-r10 ledger and guard update}
\label{subsec:v490-r10-ledger-guard}

\[
R_{\rm core}^{v490\text{-}r10}=(0,0,1)
\]
for (a reading of the factor offered, a class of shape-(iii) witnesses closed under it, a derivation of the reading open).

The layer spends no guard:
\[
R_{\rm glob}^{(v490\text{-}r10,+)}=(0,1,1,1,1),
\qquad
\Gamma_{\rm guard}^{v490\text{-}r10}=(1,1,1,1,1),
\qquad
\beta_4=1/256\ \text{unspent}.
\]
It does not imply exact finite \(C5\) zero closure, exact all-readout terminality,
\(C6/T2\), a global photon-clock/RH theorem, or global all-readout closure.

% ============================================================================
\clearpage
\section*{v4.9.0-r11 Source-Only: A Minimal Dynamics for Reading B --- Kicked Duality Rotation}
\addcontentsline{toc}{section}{v4.9.0-r11 Source-Only: A Minimal Dynamics for Reading B --- Kicked Duality Rotation}
\label{sec:v490-r11-kicked-duality-dynamics}
% ============================================================================

\paragraph{Purpose.}
r10 found that the open threads of the core meet at one gap: Reading~B has no dynamics (r10 of v4.8.0).  The C5 factor
\(3/4\) would be derived if a dynamics produced the alternation ``rotate, then compensate'' of r10.1.  A shape-(iii) witness
would need a dynamics that changes the invariant \(s\) (r10.2).  This layer proposes the smallest dynamics with these
properties and checks what it does.  It is a model proposal, not a derivation from the source.  The script is
\texttt{llm/k52\_kick\_dynamik.py}.  No value is compared with data.  Earlier layers are not rewritten.
This layer does not assert or approach an exact finite C5 zero closure.

\noindent\textbf{Source-only status.}
This v4.9.0-r11 layer is source-only.  It does not build a PDF, does not assign a new DOI,
and does not authorize a public release.  The v4.9.0 Version DOI
\href{https://doi.org/10.5281/zenodo.23534954}{10.5281/zenodo.23534954} and the Concept DOI
\href{https://doi.org/10.5281/zenodo.19210728}{10.5281/zenodo.19210728} are unchanged by it.

\subsection*{v4.9.0-r11 the proposal}
\label{subsec:v490-r11-proposal}

\paragraph{Definition (kicked duality dynamics, proposal).}
Between kicks the field turns in the duality circle at a constant rate, \(z(t)=e^{i\omega t}z\), on the doublet
\(F\mapsto FR_{\omega t}^{\mathsf T}\).  At the times \(t_n=n\tau\) the compensator \(C_{\HC}\) is applied, and the content of the
partner port is handed to a reservoir:
\[
F_{n+1}=F_nR_{\omega\tau}^{\mathsf T}C_{\HC},\qquad
\mathcal E_{\rm res}^{(n+1)}=\mathcal E_{\rm res}^{(n)}+\|F_nR_{\omega\tau}^{\mathsf T}A\|^2 .
\]
The rotation is the transport of Reading~B, the kick is the readout of the main body, and the reservoir is the partner
sector \(\operatorname{ran}A\) of r03.1.

\subsection*{v4.9.0-r11 what it does}
\label{subsec:v490-r11-properties}

\paragraph{Theoremlet v4.9.0-r11.1.}
\begin{enumerate}[leftmargin=2em]
\item[(a)] Each stage keeps \(\cos^2\omega\tau\) of the content and exports \(\sin^2\omega\tau\).  Retained plus exported
content is conserved.
\item[(b)] The stage factor equals \(3/4\) for exactly one \(\omega\tau\in(0,\pi/2)\), namely \(\omega\tau=\pi/6\).  Then the
retained content runs through \(3/4,9/16,27/64,81/256,\ldots\) and first falls below \(1/256\) at stage \(20\).
\item[(c)] In the state disk of r04 a stage is a rotation of the Bloch vector by \(2\omega\tau=60^\circ\) followed by a
measurement along the \(J\)-axis, kept with probability \((1+\cos60^\circ)/2=3/4\).
\item[(d)] The duality phase accumulates in steps of \(30^\circ\).  Three stages make the quarter turn from matter to
tachyon, six the half turn \(z\mapsto-z\) (the charge conjugation of r02.2), twelve the full turn, twenty-four the return of
the spinor of r06.4.
\item[(e)] The rotation leaves \(|s|\) unchanged.  The kick changes it: it sets \(\operatorname{Re}s=0\) and lowers
\(|s|\) (in the checked example from \(2.15\) to \(0.83\)).  The dynamics therefore leaves the transport group of
Reading~B, as r10.2 requires of any witness that changes \(s\).
\item[(f)] Dividing the same phase \(\pi/6\) into \(m\) kicks keeps \(\cos^{2m}(\pi/6m)\): \(0.75\), \(0.87\), \(0.95\), \(0.99\)
for \(m=1,2,5,50\).  The stage factor is set by how often the readout acts per phase, which is the quantum Zeno pattern.
\end{enumerate}

\paragraph{Proof sketch.}
(a)--(b) \(C_{\HC}R_\alpha C_{\HC}=\cos\alpha\,C_{\HC}\) (r10.1) and \(C_{\HC}+A=1\).  (c) \(G\mapsto RGR^{\mathsf T}\) turns the
Bloch vector by twice the angle (r04.2).  (e) direct computation.  All items were checked numerically, (b) also in exact
rational arithmetic.  \(\square\)

\subsection*{v4.9.0-r11 what it does not do}
\label{subsec:v490-r11-limits}

\begin{enumerate}[leftmargin=2em]
\item \emph{No finite closure.}  The retained content after \(n\) stages is \((3/4)^n>0\) for every \(n\).  It vanishes in
one stage only for \(\omega\tau=\pi/2\), and then the whole content sits in the reservoir: it is exported, not annihilated.
The dynamics reproduces \(C5_{\rm det/eff}\) at \(\theta=1/256\), \(n=20\), and keeps
\(0_{\rm termdata}\not\Rightarrow0_{\rm all}\).  The all-readout slot includes the reservoir.
\item \emph{No scale.}  Only the product \(\omega\tau\) is fixed.  The source supplies neither a rate \(\omega\) nor an interval
\(\tau\).  The oscillator frequency of Theorem~\ref{thm:spring}, \(\gstring\sqrt{2/3}\), is a dimensionless candidate for
\(\omega\) in the source's units.  Without a unit of time it fixes nothing.
\item \emph{No mechanism for the kick.}  Why the readout acts at regular intervals, and three times per quarter turn, is
not explained.  The triolic division of r10 is assumed, not derived.
\item \emph{Not the field equations.}  The proposal acts on the doublet at a point.  It is not a field theory and makes no
claim about the Maxwell or Einstein limits of the main body.
\end{enumerate}

\paragraph{Remark.}
The structure, a fixed polarisation rotation followed by a polariser at every cycle, is the one used in quantum Zeno
experiments with photons, among them the interaction-free measurements of Kwiat and co-workers (1995).  There a rotation of
\(\pi/(2N)\) per cycle with \(N\) cycles per quarter turn leaves \(\cos^{2N}(\pi/2N)\).  The proposal of this layer is that
structure with \(N=3\).  The correspondence concerns the algebra only.  On a light wave the compensator of r03 is not a
polariser (r03.2).

\subsection*{v4.9.0-r11 register entries}
\label{subsec:v490-r11-register}

\begin{center}
\small
\begin{tabular}{@{}>{\raggedright\arraybackslash}p{0.34\textwidth}>{\raggedright\arraybackslash}p{0.58\textwidth}@{}}
\toprule
statement & status after v4.9.0-r11\\
\midrule
dynamics of Reading~B (r10 of v4.8.0) & proposal: kicked duality rotation with compensator kicks and partner reservoir; model decision, not derived\\
C5 factor \(3/4\) & produced by the proposal at \(\omega\tau=\pi/6\), uniquely in \((0,\pi/2)\) (r11.1)\\
shape (iii) & the kick changes \(s\), but the dynamics exports instead of annihilating; no finite closure (r11)\\
scale of the dynamics & open: only \(\omega\tau\) fixed\\
\bottomrule
\end{tabular}
\end{center}

\subsection*{v4.9.0-r11 ledger and guard update}
\label{subsec:v490-r11-ledger-guard}

\[
R_{\rm dyn}^{v490\text{-}r11}=(0,0,1)
\]
for (a minimal dynamics proposed, its stage factor and cycles established, its scale and kick mechanism open).

The layer spends no guard:
\[
R_{\rm glob}^{(v490\text{-}r11,+)}=(0,1,1,1,1),
\qquad
\Gamma_{\rm guard}^{v490\text{-}r11}=(1,1,1,1,1),
\qquad
\beta_4=1/256\ \text{unspent}.
\]
It does not imply exact finite \(C5\) zero closure, exact all-readout terminality,
\(C6/T2\), a global photon-clock/RH theorem, or global all-readout closure.

% ============================================================================
\clearpage
\section*{v4.9.0-r12 Source-Only: A Reconstruction of the Bound \texorpdfstring{\(\Delta_\gamma<1/5\)}{Delta gamma < 1/5} in Place of the Missing r65}
\addcontentsline{toc}{section}{v4.9.0-r12 Source-Only: A Reconstruction of the Bound Delta gamma < 1/5 in Place of the Missing r65}
\label{sec:v490-r12-delta-gamma-reconstruction}
% ============================================================================

\paragraph{Purpose.}
The bound \(\Delta_\gamma<1/5\) carried by \(W_{\rm clk}^{\rm loc}\) comes from r65, whose full text is not available.  The
author confirmed on 2026-10-22 that it is lost and asked for a reconstruction.  r90 obtained a bound of the right size
(\(2/9\)).  r92 obtained \(1/5\) exactly as a relative gain reserve, without stating why that reading and that chain apply.
This layer states the premises explicitly, derives the strict bound from them, and draws the consequences under the
dynamics of r11.  It replaces the missing derivation.  It does not recover the r65 text and does not claim that r65 argued
this way.  The script is \texttt{llm/k53\_r65\_rekonstruktion.py}.  Earlier layers are not rewritten.

\noindent\textbf{Source-only status.}
This v4.9.0-r12 layer is source-only.  It does not build a PDF, does not assign a new DOI,
and does not authorize a public release.  The v4.9.0 Version DOI
\href{https://doi.org/10.5281/zenodo.23534954}{10.5281/zenodo.23534954} and the Concept DOI
\href{https://doi.org/10.5281/zenodo.19210728}{10.5281/zenodo.19210728} are unchanged by it.

\subsection*{v4.9.0-r12 premises}
\label{subsec:v490-r12-premises}

\begin{enumerate}[leftmargin=2em]
\item[(P1)] \emph{Chart data} of r14 (\S\ref{prop:v470-r14-finite-flowrate-classification}): overlap flux
\(S=N+\kappa N(N-1)\), relay score \(H=GS\), relay threshold \(16\).
\item[(P2)] \emph{Survival} of a window is the strict relay condition of
Definition~\ref{def:v470-r23-boundary-source-stability-window}, applied to the perturbed gain:
\(G(1-\Delta_\gamma)\,S>16\).
\item[(P3)] \emph{The perturbation is a multiplicative loss.}  \(\Delta_\gamma\) is the fraction of the clock content lost per
readout.  Under the dynamics of r11 this is the fraction exported to the partner reservoir, \(\sin^2\alpha\) for a duality
drift \(\alpha\) between readouts.
\item[(P4)] \emph{Scope.}  The guard concerns the charts that carry clock windows (unassisted score below \(16\), assisted
score at least \(16\)) on the primary relay chain \(\mathcal R_A\) of
Section~\ref{sec:v470-r15-zero-clock-phase-window-relay-stitching}, that is \(V3\) and \(V5\).
\end{enumerate}
(P1) and (P2) are in the source.  (P3) chooses r92's reading over r90's, and gives it a mechanism.  (P4) is r92's
restriction, consistent with the local index of \(W_{\rm clk}^{\rm loc}\).

\subsection*{v4.9.0-r12 the bound}
\label{subsec:v490-r12-bound}

\paragraph{Theoremlet v4.9.0-r12.1.}
Under (P1)--(P4), \(W_{\rm clk}^{\rm loc}\) is preserved if and only if
\[
\Delta_\gamma<\min_{V\in\{V3,V5\}}\Bigl(1-\frac{16}{G_VS_V}\Bigr)=\min\Bigl\{\frac15,\frac9{25}\Bigr\}=\frac15 .
\]
The inequality is strict.  At \(\Delta_\gamma=1/5\) chart \(V3\) has \(G(1-\Delta_\gamma)S=16\) exactly, which the strict window
excludes.

\paragraph{Proof sketch.}
(P2) rearranges to \(\Delta_\gamma<1-16/(GS)\).  \(V3\): \(N=5\), \(\kappa=1/4\), \(G=2\), \(S=10\), \(H=20\), hence \(1/5\).
\(V5\): \(N=4\), \(\kappa=1/2\), \(G=5/2\), \(S=10\), \(H=25\), hence \(9/25\).  Checked in exact rational arithmetic.  \(\square\)

\paragraph{Where the five comes from.}
\(V3\) is the only chart of the r14 family with \(\kappa(N-1)=1\): its pair-overlap flux \(\kappa N(N-1)\) equals its direct
flux \(N\), so \(S=2N\).  With \(G=2\) the bound becomes \(1-16/(4N)=1-4/N\), which is \(1/5\) at \(N=5\).  The binding chart is the
balanced chart of the primary chain.

\subsection*{v4.9.0-r12 consequences under the dynamics of r11}
\label{subsec:v490-r12-consequences}

\paragraph{Theoremlet v4.9.0-r12.2 (conditional on (P3)).}
\begin{enumerate}[leftmargin=2em]
\item[(a)] The clock selector survives a duality drift \(\alpha\) per readout exactly when \(\sin^2\alpha<1/5\), that is
\(\tan\alpha<1/2\), \(\alpha<26.57^\circ\).
\item[(b)] The C5 stage of r10 and r11 (\(\alpha=\pi/6\)) exports \(1/4>1/5\) and lies outside.  The tolerated collar row of
r05 (ratio \(5/6\)) exports \(1/6<1/5\) and lies inside.
\item[(c)] The rows of r05 that reach \(\Phi_5=0\) in one step, \(Z_5^{\rm mix\text{-}cancel}\) and
\(Z_5^{\rm all\text{-}cancel}\), move the whole score \(-3/4\) into \(A_4\) or \(U_6\) and are blocked by \(C4\) and \(C6\).  This is
the quarter turn of r10.2 in ledger form: the compensator port is emptied by relocating the content to another slot.
\end{enumerate}

\paragraph{Proof sketch.}
(a) (P3) in Theoremlet r12.1.  (b) \(1-3/4\) and \(1-5/6\).  (c) read off
Definition~\ref{def:v470-r05-finite-c5-iteration-ledger}.  \(\square\)

\paragraph{Reading.}
The strict C5 lowering and the clock selector are separated by the bound: a C5 stage drifts further per readout than the
clock window tolerates.  This is consistent with the two being distinct guard items, \(C5_{\rm det/eff}\) and
\(W_{\rm clk}^{\rm loc}\), established on different charts.  The r05 collar row, which r05 calls ``useful as a stability
check'', is the one whose drift the clock tolerates.  That r05 classified it as tolerated for this reason is an
interpretation.

\subsection*{v4.9.0-r12 status of the bound}
\label{subsec:v490-r12-status}

\begin{center}
\small
\begin{tabular}{@{}>{\raggedright\arraybackslash}p{0.30\textwidth}>{\raggedright\arraybackslash}p{0.62\textwidth}@{}}
\toprule
layer & contribution\\
\midrule
r87 & no derivation present; not reconstructible from the window layer alone\\
r90 & coherence reading, bound \(2/9\), \(1/5\) as rounding\\
r92 & gain-reserve reading, \(1/5\) exactly on \(\mathcal R_A\), premises implicit\\
v4.9.0-r12 & premises (P1)--(P4) stated, strictness from r23, multiplicative loss from r11, balanced chart \(V3\) identified\\
\bottomrule
\end{tabular}
\end{center}
The provenance component of \(R_{\rm cal}^{r87}\) stays open, since the r65 text itself is not recovered.  The derivation
component is closed under (P1)--(P4).

\subsection*{v4.9.0-r12 register entries}
\label{subsec:v490-r12-register}

\begin{center}
\small
\begin{tabular}{@{}>{\raggedright\arraybackslash}p{0.34\textwidth}>{\raggedright\arraybackslash}p{0.58\textwidth}@{}}
\toprule
statement & status after v4.9.0-r12\\
\midrule
\(\Delta_\gamma<1/5\) in \(W_{\rm clk}^{\rm loc}\) & derived under (P1)--(P4), strict (r12.1); r65 text lost (author), provenance open\\
binding chart & \(V3\), the balanced chart \(\kappa(N-1)=1\) of \(\mathcal R_A\)\\
clock vs C5 stage & tolerated drift \(\tan\alpha<1/2\); C5 stage outside, r05 collar row inside (r12.2)\\
r05 zero rows & relocation into \(A_4\), \(U_6\), blocked by \(C4\), \(C6\); same structure as r10.2\\
\bottomrule
\end{tabular}
\end{center}

\subsection*{v4.9.0-r12 ledger and guard update}
\label{subsec:v490-r12-ledger-guard}

\[
R_{\rm cal}^{v490\text{-}r12}=(0,0,1)
\]
for (derivation under stated premises closed, consequences under r11 recorded, r65 provenance open).

The layer spends no guard.  \(W_{\rm clk}^{\rm loc}\) under \(\Delta_\gamma<1/5\) remains the local, permitted statement:
\[
R_{\rm glob}^{(v490\text{-}r12,+)}=(0,1,1,1,1),
\qquad
\Gamma_{\rm guard}^{v490\text{-}r12}=(1,1,1,1,1),
\qquad
\beta_4=1/256\ \text{unspent}.
\]
It does not imply exact finite \(C5\) zero closure, exact all-readout terminality,
\(C6/T2\), a global photon-clock/RH theorem, or global all-readout closure.

% ============================================================================
\clearpage
\section*{v4.9.0-r13 Source-Only: The Time Scale of the Kicked Duality Dynamics}
\addcontentsline{toc}{section}{v4.9.0-r13 Source-Only: The Time Scale of the Kicked Duality Dynamics}
\label{sec:v490-r13-time-scale}
% ============================================================================

\paragraph{Purpose.}
The dynamics proposed in r11 fixes only the product \(\omega\tau=\pi/6\).  This layer derives the time scale from the clock
model the source already contains (Section~\ref{sec:v470-r15-zero-clock-phase-window-relay-stitching}).  It follows the
source in reading the photon-clock probe as a diagnostic timing function, not as an identity between Riemann zeros and
photons.  All statements are local, in the clock's own units, and make no global photon-clock statement.  The script is
\texttt{llm/k54\_zeitskala.py}.  Earlier layers are not rewritten.

\noindent\textbf{Source-only status.}
This v4.9.0-r13 layer is source-only.  It does not build a PDF, does not assign a new DOI,
and does not authorize a public release.  The v4.9.0 Version DOI
\href{https://doi.org/10.5281/zenodo.23534954}{10.5281/zenodo.23534954} and the Concept DOI
\href{https://doi.org/10.5281/zenodo.19210728}{10.5281/zenodo.19210728} are unchanged by it.

\subsection*{v4.9.0-r13 the rate from the tick profile}
\label{subsec:v490-r13-rate}

r15 writes the local clock tick as \(G_V(\tau)=1+(G_V^{\rm avail}-1)\cos^2\bigl((\tau-\tau_V)/\sigma_V\bigr)\).  Between two kicks of
r11 the energy in the compensator port is \(\cos^2(\omega t)\).  Both are \(\cos^2\) of an angle linear in time.

\paragraph{Theoremlet v4.9.0-r13.1.}
\begin{enumerate}[leftmargin=2em]
\item[(a)] Identifying the two profiles gives \(\omega=1/\sigma_V\): the tick width of r15 is the time the duality rotation of r11
needs for one radian.  The kick interval is \(\tau_{\rm k}=\tfrac\pi6\,\sigma_V\).
\item[(b)] The smallest probe with a profile has two terms.  For
\(\mathcal C_\gamma\propto|e^{i\gamma_a\tau}+e^{i\gamma_b\tau}|^2\) the coherent part is exactly
\(\cos^2\bigl(\tfrac12(\gamma_b-\gamma_a)\tau\bigr)\), so \(\sigma=2/(\gamma_b-\gamma_a)\).  The beat period is
\(T_{\rm b}=2\pi/(\gamma_b-\gamma_a)\).
\item[(c)] Hence \(\tau_{\rm k}=T_{\rm b}/6\) for every pair of clock frequencies.  Three stages, the quarter turn from matter to
tachyon, take half a beat.  Six stages, the half turn \(z\mapsto-z\), take one beat.  Twelve take two beats, and twenty-four,
the return of the spinor, take four.  The detector stage \(n=20\) lies at \(10/3\) beats.
\item[(d)] For the two lowest frequencies the source uses, \(\gamma_1=14.1347\ldots\) and \(\gamma_2=21.0220\ldots\):
\(\sigma=0.29039\), \(T_{\rm b}=0.91228\), \(\tau_{\rm k}=0.15205\) in clock units.
\end{enumerate}

\paragraph{Proof sketch.}
(b) \(|e^{ia}+e^{ib}|^2=4\cos^2\tfrac{b-a}2\); checked numerically to \(10^{-15}\).  (c) \(\tau_{\rm k}=\tfrac\pi6\cdot\tfrac2{\Delta\gamma}=
\tfrac16\cdot\tfrac{2\pi}{\Delta\gamma}\).  \(\square\)

\paragraph{What is fixed and what is not.}
The relative statement (c) holds for any pair of frequencies.  The absolute value (d) depends on which pair carries the
tick, and the source names none.  The first two are the minimal choice.  The clock unit itself has no conversion to
seconds in the source.  Supplying one would require identifying clock frequencies with physical frequencies.  The source
expressly declines that identification, and as a global statement it would touch the forbidden global photon-clock
statement.  The time scale is therefore derived up to the unit of the clock.

\subsection*{v4.9.0-r13 the windows of the primary chain}
\label{subsec:v490-r13-windows}

\paragraph{Theoremlet v4.9.0-r13.2.}
In the angle of r13.1(a), the half-widths \(\delta_V/\sigma_V\) of r15 are \(39.23^\circ\) for \(V3\), \(50.77^\circ\) for \(V5\)
and \(28.13^\circ\) for \(V8\).  A timing offset of one C5 stage, \(30^\circ\), stays inside every window of the primary chain
\(\mathcal R_A=\{V3,V5\}\) and leaves the window of \(V8\) on the chain \(\mathcal R_B\).

\paragraph{Proof sketch.}
Conversion of the r15 table.  \(\square\)

\paragraph{Reading.}
This is a second, independent reason for the restriction (P4) of r12 to \(\mathcal R_A\): only there does a single C5 stage
fit into the clock window.  The tolerance \(26.57^\circ\) of r12 concerns the loss of clock content per readout, a different
quantity from a timing offset within one window, so the two numbers do not contradict each other.  The window of \(V8\)
lies between them.

\subsection*{v4.9.0-r13 the stage angle between Hopf and lattice}
\label{subsec:v490-r13-hopf-lattice}

The main body relates the Hopf value and the measured value of the offset by
\(\bsym=\tfrac1{\sqrt3}\exp\bigl(-\gamma_1/(4\pi\gL^5)\bigr)=0.5492\) (the table with ``Hopf (continuum \(S^3\))'' and
``Physical (prime lattice)'' rows).  The stage angle \(\arctan\bsym\) of r10 is therefore \(30^\circ\) in the continuum limit
\(\gamma_1\to0\) and \(28.78^\circ\) on the lattice, with \(q=0.7683\).  The C5 factor \(3/4\) is the continuum value.  On the
lattice the Hopf reading would give stage \(22\) (r10), while the triolic reading keeps \(\pi/6\).  Which of the two describes
the C5 recursion is open.  The detector stage \(n=20\) favours the continuum value or the triolic reading.

\subsection*{v4.9.0-r13 register entries}
\label{subsec:v490-r13-register}

\begin{center}
\small
\begin{tabular}{@{}>{\raggedright\arraybackslash}p{0.34\textwidth}>{\raggedright\arraybackslash}p{0.58\textwidth}@{}}
\toprule
statement & status after v4.9.0-r13\\
\midrule
time scale of r11 & \(\omega=1/\sigma_V\), \(\tau_{\rm k}=T_{\rm b}/6\) (one sixth of the clock beat) for every frequency pair (r13.1)\\
absolute value & \(\tau_{\rm k}=0.15205\) clock units for the two lowest frequencies; no conversion to seconds, by the source's own restriction\\
r12 premise (P4) & second reason: one C5 stage fits only into the windows of \(\mathcal R_A\) (r13.2)\\
C5 factor \(3/4\) & continuum value of the Hopf reading; lattice value \(0.7683\); triolic reading unaffected\\
\bottomrule
\end{tabular}
\end{center}

\subsection*{v4.9.0-r13 ledger and guard update}
\label{subsec:v490-r13-ledger-guard}

\[
R_{\rm time}^{v490\text{-}r13}=(0,0,1)
\]
for (relative time scale derived, windows of the primary chain checked, unit of the clock open).

The layer spends no guard:
\[
R_{\rm glob}^{(v490\text{-}r13,+)}=(0,1,1,1,1),
\qquad
\Gamma_{\rm guard}^{v490\text{-}r13}=(1,1,1,1,1),
\qquad
\beta_4=1/256\ \text{unspent}.
\]
It does not imply exact finite \(C5\) zero closure, exact all-readout terminality,
\(C6/T2\), a global photon-clock/RH theorem, or global all-readout closure.

% ============================================================================
\clearpage
\section*{v4.9.0-r14 Source-Only: Why Six Kicks per Beat --- Readout Directions from the Triole and the Umkehrwalze}
\addcontentsline{toc}{section}{v4.9.0-r14 Source-Only: Why Six Kicks per Beat --- Readout Directions from the Triole and the Umkehrwalze}
\label{sec:v490-r14-readout-directions}
% ============================================================================

\paragraph{Purpose.}
r13 fixed the kick interval of the dynamics of r11 at one sixth of the clock beat.  It left open why the compensator acts
six times per beat.  This layer offers a reason built from two objects of the source: the triole and the involutive
Umkehrwalze \(J_B\) of r02.3.  The reason is an interpretation (admitted by the author, 2026-10-22), not a derivation.  The
script is \texttt{llm/k55\_kick\_richtungen.py}.  Earlier layers are not rewritten.

\noindent\textbf{Source-only status.}
This v4.9.0-r14 layer is source-only.  It does not build a PDF, does not assign a new DOI,
and does not authorize a public release.  The v4.9.0 Version DOI
\href{https://doi.org/10.5281/zenodo.23534954}{10.5281/zenodo.23534954} and the Concept DOI
\href{https://doi.org/10.5281/zenodo.19210728}{10.5281/zenodo.19210728} are unchanged by it.

\subsection*{v4.9.0-r14 kicks as passages of readout directions}
\label{subsec:v490-r14-passages}

Replace the periodic kick of r11 by an event: the compensator acts whenever the phase of the invariant \(s\) passes one of a
fixed set of readout directions.  Under a uniform duality rotation the events are periodic, and their spacing is set by the
set of directions.

\paragraph{Theoremlet v4.9.0-r14.1.}
In the \(s\)-plane, \(\arg s=2\alpha\) with \(\alpha\) the duality angle.
\begin{enumerate}[leftmargin=2em]
\item[(a)] The triole directions \(\arg s\in\{0^\circ,120^\circ,240^\circ\}\) and their images under \(J_B\), which acts as
\(s\mapsto-\bar s\) (r02.3), are together the six directions \(0^\circ,60^\circ,\dots,300^\circ\).  The set contains matter
(\(0^\circ\)) and tachyon (\(180^\circ\)).
\item[(b)] Passing them under a uniform rotation gives kicks every \(30^\circ\) of duality angle, hence the stage factor
\(\cos^2 30^\circ=3/4\).  The triole alone gives \(60^\circ\) and \(q=1/4\).  A quarter grid gives \(45^\circ\) and \(q=1/2\).
Only triole and \(J_B\) together give the factor of the source.
\item[(c)] One revolution of \(s\) around the photon (\(\arg s\) from \(0^\circ\) to \(360^\circ\)) passes the six directions once.
It is the duality half turn \(z\mapsto-z\), the charge conjugation of r02.2, and by r13 it lasts exactly one beat.
\end{enumerate}

\paragraph{Proof sketch.}
\(-\overline{e^{i\theta}}=e^{i(\pi-\theta)}\) maps \(\{0,120,240\}\) to \(\{180,60,300\}\).  The rest follows from
\(\arg s=2\alpha\) and r10.1.  Checked numerically.  \(\square\)

\paragraph{Reading.}
The six kicks per beat are then not an external schedule but the six readout directions met by one revolution of \(s\)
around the photon.  The factor \(3/4\), the detector stage \(20\), the beat, the quarter turn from matter to tachyon and the
charge conjugation of the round trip become aspects of one construction: a uniform duality rotation read out at
\(\text{triole}\cup J_B(\text{triole})\).

\paragraph{What is assumed.}
The triole of the source (Theoremlet r06.2) is a condition on the three components of a vector, \((1,\omega,\omega^2)\), not on
the phase of \(s\).  Placing the triole on \(\arg s\) is the interpretive step of this layer.  Squaring maps the triole of
phases onto a triole of phases, so the placement is consistent, but the source does not make it.  Why the readout should
act at discrete directions at all is also assumed.

\paragraph{Remark, offered and not adopted.}
Six directions at \(60^\circ\) with the triole permutations as symmetry are the root system \(A_2\) of \(SU(3)\), whose Weyl
group is \(S_3\).  r15 of v4.8.0 met \(SU(3)\) as the stabiliser of the symmetric point, there read as colour.  Whether the
readout hexagon and the colour \(SU(3)\) are related would need a map from the duality plane to a Cartan plane, which the
source does not supply.

\subsection*{v4.9.0-r14 register entries}
\label{subsec:v490-r14-register}

\begin{center}
\small
\begin{tabular}{@{}>{\raggedright\arraybackslash}p{0.34\textwidth}>{\raggedright\arraybackslash}p{0.58\textwidth}@{}}
\toprule
statement & status after v4.9.0-r14\\
\midrule
six kicks per beat (r13) & reading offered: readout at \(\text{triole}\cup J_B(\text{triole})\) in the \(s\)-plane (r14.1)\\
factor \(3/4\) & needs both triole and \(J_B\); either alone, or a quarter grid, gives another factor\\
beat & one revolution of \(s\) around the photon, equal to the charge conjugation of r02.2\\
\(A_2\)/colour & remark only\\
\bottomrule
\end{tabular}
\end{center}

\subsection*{v4.9.0-r14 ledger and guard update}
\label{subsec:v490-r14-ledger-guard}

\[
R_{\rm kick}^{v490\text{-}r14}=(0,0,1)
\]
for (a reason for the kick schedule offered, its consequences consistent with r10--r13, the placement of the triole open).

The layer spends no guard:
\[
R_{\rm glob}^{(v490\text{-}r14,+)}=(0,1,1,1,1),
\qquad
\Gamma_{\rm guard}^{v490\text{-}r14}=(1,1,1,1,1),
\qquad
\beta_4=1/256\ \text{unspent}.
\]
It does not imply exact finite \(C5\) zero closure, exact all-readout terminality,
\(C6/T2\), a global photon-clock/RH theorem, or global all-readout closure.

% ============================================================================
\clearpage
\section*{v4.9.0-r15 Source-Only: Consolidated Register of v4.9.0-r01 to r14 and the Ways Out of r17 after r08}
\addcontentsline{toc}{section}{v4.9.0-r15 Source-Only: Consolidated Register of v4.9.0-r01 to r14 and the Ways Out of r17 after r08}
\label{sec:v490-r15-consolidated-register}
% ============================================================================

\paragraph{Purpose.}
Fourteen layers since the 4.9.0 release have changed the status of several statements of the main body and added
readings of their own.  This layer collects the statuses in one register, in the manner of
\S\ref{sec:v480-r09-p2-p3-reclassification}, and lists the assumptions that only the author can confirm.  It also draws one
consequence that r08 makes available: r17 left three ways out for the cosmological fractions, and r08 decides between
them.  No new computation and no new prediction enter.  Earlier layers are not rewritten.

\noindent\textbf{Source-only status.}
This v4.9.0-r15 layer is source-only.  It does not build a PDF, does not assign a new DOI,
and does not authorize a public release.  The v4.9.0 Version DOI
\href{https://doi.org/10.5281/zenodo.23534954}{10.5281/zenodo.23534954} and the Concept DOI
\href{https://doi.org/10.5281/zenodo.19210728}{10.5281/zenodo.19210728} are unchanged by it.

\subsection*{v4.9.0-r15 the ways out of r17 after r08}
\label{subsec:v490-r15-r17-ways}

r17 found that Theorem~\ref{thm:dark-energy}, held at all epochs, needs \(\bsym\) to decrease from about \(1\) at
recombination to \(0.558\) today.  Propositions~\ref{prop:H0-tension} and~\ref{prop:saturation} assume the opposite.  It
recorded three ways out: (a) the theorem holds only at the present epoch; (b) \(\bsym\) runs downward; (c) the cascade
probabilities are some other function of the fractions at other epochs.

\paragraph{Theoremlet v4.9.0-r15.1.}
Under the rules of r08 (photon-frame values and \(k=0\) fixed, matter-frame values drifting with the local factor,
\(k(\alpha)=1\)):
\begin{enumerate}[leftmargin=2em]
\item[(b)] is excluded.  A decrease of \(\bsym\) from about \(1\) to \(0.558\) changes the local factor by far more than the
\(5\%\) that r08.1 already finds incompatible with atomic clocks.
\item[(a)] and (c) remain.  Both treat \(\bsym\) as fixed and read Theorem~\ref{thm:dark-energy} at one epoch, (a) directly,
(c) through a time-dependent bridge with fixed \(\bsym\) that the source does not supply.
\item Item~(1) of Proposition~\ref{prop:saturation} lets the neutrino mass grow with \(\gL(t)^4\).  By the clock bound on the
same local factor this growth is excluded too.  Together with item~(2), which r17 already rejected, the proposition loses
its basis.
\end{enumerate}

\paragraph{Proof sketch.}
(b): a local factor diverging towards recombination; the clock bound of r08.1 applies a fortiori.  (a), (c): neither needs a
drift.  Item~(1): \(k(m_\nu)=4\) drifts faster than \(\alpha\) with the same factor.  \(\square\)

\paragraph{Consequence.}
With r08, Theorem~\ref{thm:dark-energy} is a statement about the present epoch, \(1-\bsym^2=0.698\) against \(0.689\).
Its cosmological reach beyond that epoch is not supported by any statement of the source that survives the clocks.

\subsection*{v4.9.0-r15 register}
\label{subsec:v490-r15-register}

\begin{center}
\small
\begin{tabular}{@{}>{\raggedright\arraybackslash}p{0.30\textwidth}>{\raggedright\arraybackslash}p{0.14\textwidth}>{\raggedright\arraybackslash}p{0.48\textwidth}@{}}
\toprule
statement & layer & status\\
\midrule
conversion law, \S\ref{sec:phase-constants} & r01 & factor \(\gL\) is matter--tachyon; ways (A)--(C) open; on-shell \(\cot\theta_W\) lies at (A) (r07, observation)\\
Thm.~\ref{thm:tribonacci} & r01 & equivalent to \(\cosh2\eta=e^\eta\)\\
Prop.~\ref{prop:hawking-gBI-tachyon} (ii) & r01 & frame order does not match the axis\\
Umkehrwalze & r02 & round trip with side change \(=C=PT\); involutive form antilinear, \(J_B\colon E\leftrightarrow B\)\\
Thm.~\ref{thm:core}, compensator & r02, r03 & output \(\operatorname{Re}s=0\), not null; for a charge a dyon\\
Prop.~\ref{prop:compensator} & r03 & orthogonality holds for the product \(z\times J_Bz\), not the average\\
compensator as matrix & r04 & \(\tfrac12(1+\sigma_x)\) on the doublet; state disk, light at the centre\\
Millennium matrix & r05 & on the doublet only port weighting; reading name ``layer-transport matrix''\\
photon from matter \(+\) tachyon & r05 & only for crossed orientation\\
tachyon hit chain & r06 & cone \(56.68^\circ\), simultaneous hits in the matter frame\\
electroweak mass matrix & r06 & port weighting with massless photon port (formal)\\
H14 & r07 & refuted as exact relation, undecidable at tree level\\
H15 & r07 & pre-registered, \(m_W=80.3615\pm0.0018\) GeV, comparison pending\\
Prop.~\ref{prop:H0-tension} & r08, r09 & no working mechanism; drift excluded by clocks\\
Thm.~\ref{thm:k-values}, \(\alpha\) row & r08 & \(k=1\) belongs to \(\alpha\) (closes r03 F2)\\
P5 & r09 & corrected to \(k(\alpha)/k(G_N)=1/2\)\\
local frame factor & r09 & only a weak potential coupling survives; no field law\\
C5 factor \(3/4\) & r10, r11, r14 & produced by kicked duality rotation read out at triole \(\cup\,J_B\)(triole); interpretation\\
C5 shape (iii) & r10 & witnesses from the transport group of Reading~B excluded under the reading\\
\(\Delta_\gamma<1/5\) & r12 & derived under premises (P1)--(P4); r65 text lost, provenance open\\
time scale & r13 & kick interval one sixth of the clock beat; unit of the clock open\\
Thm.~\ref{thm:dark-energy} & r15 & present-epoch statement; r17 way (b) excluded (r15.1)\\
Prop.~\ref{prop:saturation} & r15 & no basis left (items (1) and (2))\\
\bottomrule
\end{tabular}
\end{center}

\subsection*{v4.9.0-r15 assumptions for the author}
\label{subsec:v490-r15-author}

The following steps are interpretations of these layers.  Only the author can say whether they match the original
intention.
\begin{enumerate}[leftmargin=2em]
\item Theorem~\ref{thm:frame-offset} speaks about the fixed structural factor, not the local one (r09).  This resolves its two
conflicts with Theorem~\ref{thm:k-values}.
\item The C5 recursion is an alternation of duality rotation and compensation (r10, r11).
\item The readout acts at discrete directions, and the triole lies on the phase of \(s\) (r14).
\item The guard \(W_{\rm clk}^{\rm loc}\) refers to the windows of the primary chain \(\mathcal R_A\) (r12, r13).
\item Which of the ways (A)--(C) of r01 the conversion of constants takes.
\item Whether Theorem~\ref{thm:dark-energy} is meant for the present epoch only (way (a)) or through a bridge (way (c)).
\end{enumerate}

\subsection*{v4.9.0-r15 ledger and guard update}
\label{subsec:v490-r15-ledger-guard}

\[
R_{\rm reg}^{v490\text{-}r15}=(0,0,1)
\]
for (register consolidated, r17 way (b) excluded, author decisions listed and open).

The layer spends no guard:
\[
R_{\rm glob}^{(v490\text{-}r15,+)}=(0,1,1,1,1),
\qquad
\Gamma_{\rm guard}^{v490\text{-}r15}=(1,1,1,1,1),
\qquad
\beta_4=1/256\ \text{unspent}.
\]
It does not imply exact finite \(C5\) zero closure, exact all-readout terminality,
\(C6/T2\), a global photon-clock/RH theorem, or global all-readout closure.

% ============================================================================
\clearpage
\section*{v4.9.0-r16 Source-Only: Author Decisions on the Open Readings of r15}
\addcontentsline{toc}{section}{v4.9.0-r16 Source-Only: Author Decisions on the Open Readings of r15}
\label{sec:v490-r16-author-decisions}
% ============================================================================

\paragraph{Purpose.}
r15 listed six readings that only the author can confirm (\S\ref{subsec:v490-r15-author}).  The author answered four of them
on 2026-10-22.  This layer records the answers and what changes in the register.  No computation enters.  Earlier layers are
not rewritten.

\noindent\textbf{Source-only status.}
This v4.9.0-r16 layer is source-only.  It does not build a PDF, does not assign a new DOI,
and does not authorize a public release.  The v4.9.0 Version DOI
\href{https://doi.org/10.5281/zenodo.23534954}{10.5281/zenodo.23534954} and the Concept DOI
\href{https://doi.org/10.5281/zenodo.19210728}{10.5281/zenodo.19210728} are unchanged by it.

\subsection*{v4.9.0-r16 the decisions}
\label{subsec:v490-r16-decisions}

\begin{center}
\small
\begin{tabular}{@{}>{\raggedright\arraybackslash}p{0.06\textwidth}>{\raggedright\arraybackslash}p{0.40\textwidth}>{\raggedright\arraybackslash}p{0.46\textwidth}@{}}
\toprule
no. & question (r15) & author's answer (2026-10-22)\\
\midrule
1 & Does Theorem~\ref{thm:frame-offset} speak about the fixed structural factor? & \textbf{Yes, structural.}\\
2, 3 & Is the C5 recursion an alternation of duality rotation and compensation, read out at the triole and its image under \(J_B\)? & \textbf{Yes, as originally intended.}\\
5 & Which of the ways (A)--(C) does the conversion of constants take? & left open\\
6 & Is Theorem~\ref{thm:dark-energy} meant for the present epoch only, or through a bridge? & left open\\
4 & Does \(W_{\rm clk}^{\rm loc}\) refer to the windows of \(\mathcal R_A\)? & not asked; stays a premise of r12\\
\bottomrule
\end{tabular}
\end{center}

\subsection*{v4.9.0-r16 consequences}
\label{subsec:v490-r16-consequences}

\begin{enumerate}[leftmargin=2em]
\item \emph{Decision 1.}  The two conflicts between Theorem~\ref{thm:frame-offset} and Theorem~\ref{thm:k-values} recorded in
r09 are resolved by the author.  The theorem concerns \(\gL=\gBI/\gstring\) as a fixed structural ratio.  It makes no
statement about how a measured constant depends on the observer.  Any variation runs through the local factor of r08 and
r09.  Theorem~\ref{thm:k-values} stands as written, with the \(\alpha\) row read as in r08.
\item \emph{Decisions 2 and 3.}  The reading of r10, r11 and r14 is the author's intended reading of the C5 recursion: a
uniform duality rotation, compensated whenever the phase of \(s\) passes the triole or its \(J_B\)-image.  Its status changes
from ``interpretation offered'' to ``author's reading''.  It remains a reading, not a derivation from the earlier text.
The source never states the recursion in this form, and the placement of the triole on \(\arg s\) is not in r06.
Theoremlet r10.2 on shape-(iii) witnesses, which was conditional on the reading, now holds under the author's reading.
\item \emph{Decisions 5 and 6} stay open as recorded in r01 and r15.
\end{enumerate}

\subsection*{v4.9.0-r16 register entries}
\label{subsec:v490-r16-register}

\begin{center}
\small
\begin{tabular}{@{}>{\raggedright\arraybackslash}p{0.34\textwidth}>{\raggedright\arraybackslash}p{0.58\textwidth}@{}}
\toprule
statement & status after v4.9.0-r16\\
\midrule
Thm.~\ref{thm:frame-offset} & structural statement about \(\gL\) (author); conflicts of r09 resolved\\
C5 recursion (r10, r11, r14) & author's reading: kicked duality rotation read out at triole \(\cup\,J_B\)(triole)\\
C5 shape (iii) (r10.2) & witnesses from the transport group of Reading~B excluded under the author's reading\\
conversion ways (A)--(C) & open (author)\\
Thm.~\ref{thm:dark-energy}, epoch & open between r17 ways (a) and (c) (author)\\
\bottomrule
\end{tabular}
\end{center}

\subsection*{v4.9.0-r16 ledger and guard update}
\label{subsec:v490-r16-ledger-guard}

\[
R_{\rm auth}^{v490\text{-}r16}=(0,0,1)
\]
for (decisions 1 to 3 recorded, their consequences entered, decisions 5 and 6 open).

The layer spends no guard:
\[
R_{\rm glob}^{(v490\text{-}r16,+)}=(0,1,1,1,1),
\qquad
\Gamma_{\rm guard}^{v490\text{-}r16}=(1,1,1,1,1),
\qquad
\beta_4=1/256\ \text{unspent}.
\]
It does not imply exact finite \(C5\) zero closure, exact all-readout terminality,
\(C6/T2\), a global photon-clock/RH theorem, or global all-readout closure.

